\documentclass[12pt,a4paper]{article}

% --- Pacotes Essenciais ---
\usepackage[utf8]{inputenc}
\usepackage[T1]{fontenc}
% Nota: substitua 'english' por 'portuges' se a sua distribuicao LaTeX
% incluir o suporte babel para portugues (pacote texlive-lang-portuguese).
\usepackage[english]{babel}
\usepackage{amsmath}
\usepackage{amssymb}
\usepackage{amsthm}
\usepackage{geometry}
\usepackage{enumitem}
\usepackage{booktabs}
\usepackage{xcolor}
\usepackage{tcolorbox}
\usepackage{parskip}
\usepackage{titlesec}
\usepackage{fancyhdr}
\usepackage{mdframed}

\geometry{a4paper, top=2.5cm, bottom=2.5cm, left=2.5cm, right=2.5cm, headheight=15pt}

% --- Cabeçalho e rodapé ---
\pagestyle{fancy}
\fancyhf{}
\rhead{\textit{Resoluções Detalhadas}}
\lhead{\textit{Teoria dos Conjuntos, Relações e Funções}}
\cfoot{\thepage}
\renewcommand{\headrulewidth}{0.4pt}

% --- Formatação de secções ---
\titleformat{\section}[block]{\large\bfseries\color{blue!70!black}}{}{0em}{}[\titlerule]
\titleformat{\subsection}[runin]{\bfseries}{}{0em}{}[.\enspace]

% --- Ambiente para conceitos teóricos ---
\newmdenv[
  backgroundcolor=blue!5,
  linecolor=blue!40!black,
  linewidth=1pt,
  roundcorner=4pt,
  innertopmargin=6pt,
  innerbottommargin=6pt,
  innerleftmargin=10pt,
  innerrightmargin=10pt
]{teoribox}

% --- Ambiente para notas ---
\newmdenv[
  backgroundcolor=orange!8,
  linecolor=orange!60!black,
  linewidth=1pt,
  roundcorner=4pt,
  innertopmargin=6pt,
  innerbottommargin=6pt,
  innerleftmargin=10pt,
  innerrightmargin=10pt
]{notabox}

% --- Pacotes e definições adicionais vindos de main 2.tex ---
\usepackage{array}
\usepackage[colorlinks=false, pdfborder={0 0 0}]{hyperref}

\definecolor{azul}{RGB}{0,70,140}
\definecolor{verde}{RGB}{0,100,60}
\definecolor{fundoteoria}{RGB}{240,248,255}
\definecolor{fundoex}{RGB}{250,250,240}

\newmdenv[backgroundcolor=fundoteoria,linecolor=azul,linewidth=1.2pt,
          leftmargin=0pt,rightmargin=0pt,innerleftmargin=10pt,
          innerrightmargin=10pt,innertopmargin=8pt,innerbottommargin=8pt]{teoriacaixa}

\newmdenv[backgroundcolor=fundoex,linecolor=verde,linewidth=1pt,
          leftmargin=0pt,rightmargin=0pt,innerleftmargin=10pt,
          innerrightmargin=10pt,innertopmargin=6pt,innerbottommargin=6pt]{exemplocaixa}

\newcommand{\keyword}[1]{\textbf{\textcolor{azul}{#1}}}
\newcommand{\passo}[1]{\medskip\noindent\textcolor{verde}{\textbf{Passo #1.}}\;}


% --- Título combinado ---
\title{%
  \vspace{-1cm}
  {\Large \textbf{Matemática Discreta}}\\[0.4em]
  {\large Resoluções e Explicações Detalhadas}\\[0.3em]
  {\normalsize Capítulos 1, 2, 3, 10, 12, 13, 14, 15 e Apêndice B}
}
\author{}
\date{}

\begin{document}

\maketitle
\thispagestyle{fancy}

\tableofcontents
\newpage

% =============================================================
\section*{Capítulo 1 --- Teoria dos Conjuntos}
\addcontentsline{toc}{section}{Capítulo 1 --- Teoria dos Conjuntos}
% =============================================================

% --------------------------------------------------
\section*{Exercício 1.1}
\addcontentsline{toc}{subsection}{Exercício 1.1}
% --------------------------------------------------

\textbf{Enunciado:} Quais dos seguintes conjuntos são iguais?
$\{x, y, z\}$, $\{z, y, z, x\}$, $\{y, x, y, z\}$, $\{y, z, x, y\}$.

\begin{teoribox}
\textbf{Conceito fundamental:} Dois conjuntos são \emph{iguais} se e só se contêm exactamente os mesmos elementos. Por definição, a \textbf{ordem} em que os elementos são listados é irrelevante e a \textbf{repetição} de elementos não altera o conjunto.
\end{teoribox}

\textbf{Resolução detalhada:}

Analisemos cada conjunto, eliminando repetições e ignorando a ordem:
\begin{itemize}
  \item $\{x, y, z\}$ --- elementos distintos: $x$, $y$, $z$.
  \item $\{z, y, z, x\}$ --- o elemento $z$ repete-se. Eliminando a repetição: $\{z, y, x\} = \{x, y, z\}$.
  \item $\{y, x, y, z\}$ --- o elemento $y$ repete-se. Eliminando: $\{y, x, z\} = \{x, y, z\}$.
  \item $\{y, z, x, y\}$ --- o elemento $y$ repete-se. Eliminando: $\{y, z, x\} = \{x, y, z\}$.
\end{itemize}

Após normalizar os quatro conjuntos, todos resultam em $\{x, y, z\}$. Portanto, \textbf{todos os quatro conjuntos são iguais}.

% --------------------------------------------------
\section*{Exercício 1.2}
\addcontentsline{toc}{subsection}{Exercício 1.2}
% --------------------------------------------------

\textbf{Enunciado:} Liste os elementos de cada conjunto, onde $\mathbb{N} = \{1, 2, 3, \ldots\}$.

\begin{teoribox}
\textbf{Notação de construção de conjuntos (set-builder notation):} $\{x \in A \mid P(x)\}$ denota o conjunto de todos os $x$ em $A$ que satisfazem a propriedade $P(x)$. Este método de definir conjuntos é chamado de \emph{compreensão}.
\end{teoribox}

\textbf{(a)} $A = \{x \in \mathbb{N} \mid 3 < x < 9\}$

A condição pede os inteiros positivos estritamente maiores que 3 e estritamente menores que 9. Os inteiros nesse intervalo aberto $(3,9)$ são $4, 5, 6, 7, 8$. Logo:
\[
A = \{4, 5, 6, 7, 8\}
\]

\textbf{(b)} $B = \{x \in \mathbb{N} \mid x \text{ é par},\, x < 11\}$

Procuramos os inteiros positivos pares menores que 11. Os números pares em $\mathbb{N}$ são $2, 4, 6, 8, 10, 12, \ldots$ e os que satisfazem $x < 11$ são os primeiros cinco. Logo:
\[
B = \{2, 4, 6, 8, 10\}
\]

\textbf{(c)} $C = \{x \in \mathbb{N} \mid 4 + x = 3\}$

A equação $4 + x = 3$ implica $x = -1$. Como $-1 \notin \mathbb{N}$ (os inteiros positivos começam em 1), nenhum elemento de $\mathbb{N}$ satisfaz esta condição. Logo:
\[
C = \emptyset \quad \text{(o conjunto vazio)}
\]

% --------------------------------------------------
\section*{Exercício 1.3}
\addcontentsline{toc}{subsection}{Exercício 1.3}
% --------------------------------------------------

\textbf{Enunciado:} Seja $A = \{2, 3, 4, 5\}$.

\begin{teoribox}
\textbf{Subconjunto ($\subseteq$) e subconjunto próprio ($\subsetneq$):}
\begin{itemize}
  \item $A \subseteq B$ significa que \emph{todo} elemento de $A$ também pertence a $B$.
  \item Para mostrar que $A \not\subseteq B$, basta encontrar \emph{um único} elemento $a \in A$ tal que $a \notin B$.
  \item $A \subsetneq B$ (subconjunto próprio) significa $A \subseteq B$ e $A \neq B$, ou seja, existe pelo menos um elemento em $B$ que não está em $A$.
\end{itemize}
\end{teoribox}

\textbf{(a)} Mostrar que $A$ não é subconjunto de $B = \{x \in \mathbb{N} \mid x \text{ é par}\}$.

Para refutar $A \subseteq B$, precisamos de encontrar um elemento de $A$ que não pertença a $B$. O conjunto $B$ contém apenas números pares. Verificamos:
\[
3 \in A \quad \text{e} \quad 3 \text{ é ímpar} \Rightarrow 3 \notin B
\]
Como existe um elemento ($3$) em $A$ que não está em $B$, concluímos $A \not\subseteq B$.

\textbf{(b)} Mostrar que $A$ é subconjunto próprio de $C = \{1, 2, 3, \ldots, 9\}$.

Primeiro verificamos $A \subseteq C$: cada elemento de $A$ é um inteiro entre 2 e 5, portanto pertence a $C = \{1,\ldots,9\}$:
\[
2, 3, 4, 5 \in C \quad \checkmark
\]
De seguida, verificamos que $A \neq C$: o elemento $1 \in C$ mas $1 \notin A$. Assim $A \subsetneq C$.

% --------------------------------------------------
\section*{Exercício 1.4}
\addcontentsline{toc}{subsection}{Exercício 1.4}
% --------------------------------------------------

\textbf{Enunciado:} Seja $U = \{1,2,\ldots,9\}$ o conjunto universal, e sejam dados os conjuntos $A, B, C, D, E, F$. Calcular uniões e intersecções.

\begin{teoribox}
\textbf{Operações de conjuntos:}
\begin{itemize}
  \item \textbf{União:} $X \cup Y = \{x \mid x \in X \text{ ou } x \in Y\}$ --- todos os elementos que estão em pelo menos um dos conjuntos.
  \item \textbf{Intersecção:} $X \cap Y = \{x \mid x \in X \text{ e } x \in Y\}$ --- apenas os elementos comuns a ambos.
  \item \textbf{Teorema 1.4:} Se $F \subseteq D$, então $D \cup F = D$ e $D \cap F = F$.
\end{itemize}
\end{teoribox}

Dados: $A = \{1,2,3,4,5\}$, $B = \{4,5,6,7\}$, $C = \{5,6,7,8,9\}$, $D = \{1,3,5,7,9\}$, $F = \{1,5,9\}$.

\textbf{(a)}
\begin{align*}
A \cup B &= \{1,2,3,4,5\} \cup \{4,5,6,7\} = \{1,2,3,4,5,6,7\} \\
A \cap B &= \{4,5\} \quad \text{(elementos comuns a } A \text{ e } B\text{)}
\end{align*}

\textbf{(b)}
\begin{align*}
A \cup C &= \{1,2,3,4,5\} \cup \{5,6,7,8,9\} = \{1,2,3,4,5,6,7,8,9\} = U \\
A \cap C &= \{5\}
\end{align*}

\textbf{(c)} Observemos que $F = \{1,5,9\} \subseteq \{1,3,5,7,9\} = D$. Pelo Teorema 1.4:
\begin{align*}
D \cup F &= D = \{1,3,5,7,9\} \\
D \cap F &= F = \{1,5,9\}
\end{align*}

% --------------------------------------------------
\section*{Exercício 1.5}
\addcontentsline{toc}{subsection}{Exercício 1.5}
% --------------------------------------------------

\textbf{Enunciado:} Considerar os conjuntos do Exercício 1.4. Calcular complementos, diferenças e diferenças simétricas.

\begin{teoribox}
\textbf{Operações adicionais:}
\begin{itemize}
  \item \textbf{Complemento:} $X^c = \{x \in U \mid x \notin X\}$ --- elementos do universo que não estão em $X$.
  \item \textbf{Diferença:} $X \setminus Y = \{x \mid x \in X \text{ e } x \notin Y\}$ --- elementos de $X$ que não pertencem a $Y$.
  \item \textbf{Diferença simétrica:} $X \oplus Y = (X \setminus Y) \cup (Y \setminus X)$ --- elementos que estão num ou noutro, mas não em ambos.
\end{itemize}
\end{teoribox}

Com $U = \{1,2,\ldots,9\}$, $A=\{1,2,3,4,5\}$, $B=\{4,5,6,7\}$, $D=\{1,3,5,7,9\}$, $E=\{2,4,6,8\}$, $F=\{1,5,9\}$.

\textbf{(a) Complementos} (elementos de $U$ que faltam em cada conjunto):
\begin{align*}
A^c &= U \setminus A = \{6,7,8,9\} \\
B^c &= U \setminus B = \{1,2,3,8,9\} \\
D^c &= U \setminus D = \{2,4,6,8\} = E \\
E^c &= U \setminus E = \{1,3,5,7,9\} = D
\end{align*}
Note a dualidade: $D^c = E$ e $E^c = D$, o que significa que $D$ e $E$ são complementares em $U$.

\textbf{(b) Diferenças}:
\begin{align*}
A \setminus B &= \{1,2,3,4,5\} \setminus \{4,5,6,7\} = \{1,2,3\} \\
B \setminus A &= \{4,5,6,7\} \setminus \{1,2,3,4,5\} = \{6,7\} \\
D \setminus E &= \{1,3,5,7,9\} \setminus \{2,4,6,8\} = \{1,3,5,7,9\} = D
\end{align*}
O último resultado confirma que $D \cap E = \emptyset$ (conjuntos disjuntos).

\textbf{(c) Diferenças simétricas}:
\begin{align*}
A \oplus B &= \{1,2,3\} \cup \{6,7\} = \{1,2,3,6,7\} \\
C \oplus D &= (C \setminus D) \cup (D \setminus C) = \{6,8\} \cup \{1,3\} = \{1,3,6,8\} \\
E \oplus F &= (E \setminus F) \cup (F \setminus E) = \{2,4,6,8\} \cup \{1,5,9\} = E \cup F
\end{align*}
O último resultado ocorre porque $E \cap F = \emptyset$, pelo que $E \oplus F = E \cup F$.

% --------------------------------------------------
\section*{Exercício 1.6}
\addcontentsline{toc}{subsection}{Exercício 1.6}
% --------------------------------------------------

\textbf{Enunciado:} Mostrar por exemplo que: (a) $A \cap B = A \cap C$ não implica $B = C$; (b) $A \cup B = A \cup C$ não implica $B = C$.

\begin{teoribox}
\textbf{Estratégia:} Para refutar uma implicação geral, basta construir um \emph{contra-exemplo} --- um caso concreto onde a hipótese é verdadeira mas a conclusão é falsa.
\end{teoribox}

\textbf{(a)} Tomemos $A = \{1,2\}$, $B = \{2,3\}$, $C = \{2,4\}$. Então:
\[
A \cap B = \{2\} = A \cap C \quad \text{mas} \quad B = \{2,3\} \neq \{2,4\} = C
\]
A intersecção com $A$ ``filtra'' apenas os elementos de $A$; o que $B$ e $C$ têm fora de $A$ é irrelevante para a intersecção, mas distingue-os.

\textbf{(b)} Tomemos $A = \{1,2\}$, $B = \{1,3\}$, $C = \{2,3\}$. Então:
\[
A \cup B = \{1,2,3\} = A \cup C \quad \text{mas} \quad B = \{1,3\} \neq \{2,3\} = C
\]
O conjunto $A$ ``cobre'' a diferença entre $B$ e $C$ na união, tornando os resultados iguais apesar de $B \neq C$.

% --------------------------------------------------
\section*{Exercício 1.7}
\addcontentsline{toc}{subsection}{Exercício 1.7}
% --------------------------------------------------

\textbf{Enunciado:} Provar que $B \setminus A = B \cap A^c$.

\begin{teoribox}
\textbf{Método:} Para provar igualdade de dois conjuntos $X = Y$, mostramos que cada elemento de $X$ pertence a $Y$ e vice-versa. Equivalentemente, usamos a lógica proposicional para traduzir as definições em condições equivalentes.
\end{teoribox}

\textbf{Demonstração} (por dupla inclusão via definições):

Um elemento $x$ pertence a $B \setminus A$ se e só se satisfaz simultaneamente:
\[
x \in B \quad \text{e} \quad x \notin A
\]
A condição $x \notin A$ é, por definição do complemento, equivalente a $x \in A^c$. Portanto:
\[
x \in B \setminus A \iff x \in B \text{ e } x \notin A \iff x \in B \text{ e } x \in A^c \iff x \in B \cap A^c
\]
Como a cadeia de equivalências é válida para qualquer $x$, conclui-se $B \setminus A = B \cap A^c$. \hfill$\square$

% --------------------------------------------------
\section*{Exercício 1.8}
\addcontentsline{toc}{subsection}{Exercício 1.8}
% --------------------------------------------------

\textbf{Enunciado:} Provar o Teorema 1.4: as seguintes afirmações são equivalentes: $A \subseteq B$,\ $A \cap B = A$,\ $A \cup B = B$.

\begin{teoribox}
\textbf{Equivalência de três afirmações:} Para provar que $P \Leftrightarrow Q \Leftrightarrow R$, basta demonstrar a cadeia $P \Rightarrow Q \Rightarrow R \Rightarrow P$ (ou qualquer outra cadeia cíclica).
\end{teoribox}

\textbf{($A \subseteq B \Rightarrow A \cap B = A$):} Suponha $A \subseteq B$. Provamos $A \cap B = A$ mostrando dupla inclusão.
\begin{itemize}
  \item $A \cap B \subseteq A$: sempre verdade por definição de intersecção.
  \item $A \subseteq A \cap B$: seja $x \in A$. Como $A \subseteq B$, temos $x \in B$. Logo $x \in A \cap B$.
\end{itemize}

\textbf{($A \cap B = A \Rightarrow A \subseteq B$):} Suponha $A \cap B = A$. Seja $x \in A$. Então $x \in A \cap B$, o que implica $x \in B$. Logo $A \subseteq B$.

\textbf{($A \subseteq B \Rightarrow A \cup B = B$):} Suponha $A \subseteq B$.
\begin{itemize}
  \item $B \subseteq A \cup B$: sempre verdade.
  \item $A \cup B \subseteq B$: seja $x \in A \cup B$. Então $x \in A$ ou $x \in B$. Se $x \in A$, então $x \in B$ pois $A \subseteq B$. Em ambos os casos $x \in B$.
\end{itemize}

\textbf{($A \cup B = B \Rightarrow A \subseteq B$):} Seja $x \in A$. Então $x \in A \cup B = B$, logo $x \in B$. \hfill$\square$

% --------------------------------------------------
\section*{Exercício 1.9}
\addcontentsline{toc}{subsection}{Exercício 1.9}
% --------------------------------------------------

\textbf{Enunciado:} Ilustrar a Lei de De Morgan $(A \cup B)^c = A^c \cap B^c$ usando diagramas de Venn.

\begin{teoribox}
\textbf{Leis de De Morgan:}
\[
(A \cup B)^c = A^c \cap B^c \qquad \text{e} \qquad (A \cap B)^c = A^c \cup B^c
\]
Estas leis estabelecem uma dualidade entre uniões e intersecções quando aplicado o complemento.
\end{teoribox}

\textbf{Raciocínio:}

O lado esquerdo, $(A \cup B)^c$, representa a região fora da reunião de $A$ e $B$, ou seja, tudo no universo $U$ que não pertence a $A$ \emph{nem} a $B$.

O lado direito, $A^c \cap B^c$, representa a intersecção das regiões exteriores a $A$ e a $B$ individualmente, ou seja, os pontos que estão simultaneamente fora de $A$ \emph{e} fora de $B$.

Descrevendo logicamente: um ponto $x$ pertence a $(A \cup B)^c$ se e só se $x \notin A \cup B$, o que significa $x \notin A$ e $x \notin B$, que é exactamente a condição para $x \in A^c \cap B^c$.

Num diagrama de Venn, ambas as expressões sombreiam a mesma região --- o exterior ao círculo que abrange $A$ e $B$ --- confirmando geometricamente a igualdade.

% --------------------------------------------------
\section*{Exercício 1.10}
\addcontentsline{toc}{subsection}{Exercício 1.10}
% --------------------------------------------------

\textbf{Enunciado:} Provar a Lei Distributiva: $A \cap (B \cup C) = (A \cap B) \cup (A \cap C)$.

\begin{teoribox}
\textbf{Analogia lógica:} Esta lei é o espelho da distributividade lógica $p \land (q \lor r) \equiv (p \land q) \lor (p \land r)$, onde $\cap$ corresponde a ``e'' ($\land$) e $\cup$ a ``ou'' ($\lor$).
\end{teoribox}

\textbf{Demonstração} (usando definições de conjunto via lógica):
\begin{align*}
A \cap (B \cup C) &= \{x \mid x \in A \text{ e } x \in B \cup C\} \\
&= \{x \mid x \in A \text{ e } (x \in B \text{ ou } x \in C)\} \\
&= \{x \mid (x \in A \text{ e } x \in B) \text{ ou } (x \in A \text{ e } x \in C)\} \\
&= (A \cap B) \cup (A \cap C)
\end{align*}
O terceiro passo usa a distributividade lógica: $p \land (q \lor r) \equiv (p \land q) \lor (p \land r)$. \hfill$\square$

% --------------------------------------------------
\section*{Exercício 1.11}
\addcontentsline{toc}{subsection}{Exercício 1.11}
% --------------------------------------------------

\textbf{Enunciado:} Escrever o dual de: (a) $(U \cap A) \cup (B \cap A) = A$;\quad (b) $(A \cap U) \cap (\emptyset \cup A^c) = \emptyset$.

\begin{teoribox}
\textbf{Princípio de Dualidade:} Numa identidade de teoria dos conjuntos, se trocarmos sistematicamente $\cup \leftrightarrow \cap$ e $U \leftrightarrow \emptyset$ (mantendo complementos), obtemos outra identidade válida --- o seu \emph{dual}.
\end{teoribox}

Aplicando a troca $\cup \leftrightarrow \cap$ e $U \leftrightarrow \emptyset$:

\textbf{(a)} $(U \cap A) \cup (B \cap A) = A$ \quad torna-se:
\[
(\emptyset \cup A) \cap (B \cup A) = A
\]

\textbf{(b)} $(A \cap U) \cap (\emptyset \cup A^c) = \emptyset$ \quad torna-se:
\[
(A \cup \emptyset) \cup (U \cap A^c) = U
\]

% --------------------------------------------------
\section*{Exercício 1.12}
\addcontentsline{toc}{subsection}{Exercício 1.12}
% --------------------------------------------------

\textbf{Enunciado:} Provar que $(A \cup B) \setminus (A \cap B) = (A \setminus B) \cup (B \setminus A)$.

\begin{teoribox}
\textbf{Diferença simétrica:} Ambos os membros da igualdade são definições equivalentes de $A \oplus B$. A prova usa $X \setminus Y = X \cap Y^c$ (Exercício 1.7) e as Leis de De Morgan.
\end{teoribox}

\textbf{Demonstração:}
\begin{align*}
(A \cup B) \setminus (A \cap B) &= (A \cup B) \cap (A \cap B)^c \\
&= (A \cup B) \cap (A^c \cup B^c) \quad \text{(De Morgan)} \\
&= \bigl[(A \cup B) \cap A^c\bigr] \cup \bigl[(A \cup B) \cap B^c\bigr] \quad \text{(Distributiva)} \\
&= \bigl[(A \cap A^c) \cup (B \cap A^c)\bigr] \cup \bigl[(A \cap B^c) \cup (B \cap B^c)\bigr] \\
&= \bigl[\emptyset \cup (B \cap A^c)\bigr] \cup \bigl[(A \cap B^c) \cup \emptyset\bigr] \\
&= (B \cap A^c) \cup (A \cap B^c) \\
&= (B \setminus A) \cup (A \setminus B) = (A \setminus B) \cup (B \setminus A)
\end{align*}
\hfill$\square$

% --------------------------------------------------
\section*{Exercício 1.13}
\addcontentsline{toc}{subsection}{Exercício 1.13}
% --------------------------------------------------

\textbf{Enunciado:} Determinar a validade do argumento:
\begin{itemize}
  \item $S_1$: Todos os meus amigos são músicos.
  \item $S_2$: John é meu amigo.
  \item $S_3$: Nenhum dos meus vizinhos é músico.
  \item Conclusão $S$: John não é meu vizinho.
\end{itemize}

\begin{teoribox}
\textbf{Raciocínio por diagramas de conjuntos:} Num diagrama de Venn, as premissas $S_1$ e $S_3$ colocam restrições geométricas. $S_1$ diz que o conjunto dos Amigos está contido no conjunto dos Músicos. $S_3$ diz que o conjunto dos Vizinhos é disjunto do conjunto dos Músicos. Pela transitividade, Amigos e Vizinhos são disjuntos.
\end{teoribox}

\textbf{Análise:}

\textit{Passo 1 ---} De $S_1$: $\text{Amigos} \subseteq \text{Músicos}$.

\textit{Passo 2 ---} De $S_3$: $\text{Vizinhos} \cap \text{Músicos} = \emptyset$.

\textit{Passo 3 ---} Das duas premissas conjuntamente: como Amigos $\subseteq$ Músicos, e Músicos $\cap$ Vizinhos $= \emptyset$, então também Amigos $\cap$ Vizinhos $= \emptyset$.

\textit{Passo 4 ---} De $S_2$: John $\in$ Amigos. Como Amigos $\cap$ Vizinhos $= \emptyset$, John $\notin$ Vizinhos.

Portanto, a conclusão $S$ é uma consequência lógica válida das premissas. O \textbf{argumento é válido}.

% --------------------------------------------------
\section*{Exercício 1.14}
\addcontentsline{toc}{subsection}{Exercício 1.14}
% --------------------------------------------------

\textbf{Enunciado:} 140 estudantes; 60 completaram $A$, 45 completaram $B$, 20 completaram ambos. Encontrar: (a) pelo menos um; (b) exactamente um; (c) nenhum.

\begin{teoribox}
\textbf{Princípio da Inclusão-Exclusão (para dois conjuntos):}
\[
n(A \cup B) = n(A) + n(B) - n(A \cap B)
\]
A subtracção de $n(A \cap B)$ corrige a dupla contagem dos elementos que foram contados tanto em $A$ como em $B$.
\end{teoribox}

\textbf{(a) Pelo menos um} ($= n(A \cup B)$):
\[
n(A \cup B) = n(A) + n(B) - n(A \cap B) = 60 + 45 - 20 = 85
\]

\textbf{(b) Exactamente um} ($= n(A \oplus B)$):

O número de estudantes que completaram $A$ mas não $B$: $60 - 20 = 40$.
O número que completaram $B$ mas não $A$: $45 - 20 = 25$. Logo:
\[
n(A \oplus B) = 40 + 25 = 65
\]

\textbf{(c) Nenhum} ($= n(A^c \cap B^c)$):
\[
n((A \cup B)^c) = n(U) - n(A \cup B) = 140 - 85 = 55
\]

% --------------------------------------------------
\section*{Exercício 1.15}
\addcontentsline{toc}{subsection}{Exercício 1.15}
% --------------------------------------------------

\textbf{Enunciado:} Num inquérito de 120 pessoas, com dados sobre leitores de Newsweek ($N$), Time ($T$) e Fortune ($F$). Encontrar: (a) lêem pelo menos uma; (b) distribuição por regiões do diagrama de Venn; (c) lêem exactamente uma.

\begin{teoribox}
\textbf{Princípio da Inclusão-Exclusão (para três conjuntos):}
\[
n(N \cup T \cup F) = n(N) + n(T) + n(F) - n(N \cap T) - n(N \cap F) - n(T \cap F) + n(N \cap T \cap F)
\]
\end{teoribox}

\textbf{(a)} Dados: $n(N)=65$, $n(T)=45$, $n(F)=42$, $n(N\cap T)=20$, $n(N\cap F)=25$, $n(T\cap F)=15$, $n(N\cap T\cap F)=8$.
\begin{align*}
n(N \cup T \cup F) &= 65 + 45 + 42 - 20 - 25 - 15 + 8 = 100
\end{align*}

\textbf{(b)} Preenchemos o diagrama de Venn de ``dentro para fora'':
\begin{align*}
\text{Só em } N \cap T \cap F &= 8 \\
\text{Em } N \cap T \text{, excluindo o centro} &= 20 - 8 = 12 \\
\text{Em } N \cap F \text{, excluindo o centro} &= 25 - 8 = 17 \\
\text{Em } T \cap F \text{, excluindo o centro} &= 15 - 8 = 7 \\
\text{Só em } N &= 65 - 12 - 8 - 17 = 28 \\
\text{Só em } T &= 45 - 12 - 8 - 7 = 18 \\
\text{Só em } F &= 42 - 17 - 8 - 7 = 10 \\
\text{Nenhuma} &= 120 - 100 = 20
\end{align*}

\textbf{(c)} Exactamente uma revista: $28 + 18 + 10 = 56$ pessoas.

% --------------------------------------------------
\section*{Exercício 1.16}
\addcontentsline{toc}{subsection}{Exercício 1.16}
% --------------------------------------------------

\textbf{Enunciado:} Provar o Teorema 1.9: Se $A$ e $B$ são finitos, então $n(A \cup B) = n(A) + n(B) - n(A \cap B)$.

\textbf{Demonstração:}

Ao contar $n(A) + n(B)$, cada elemento de $A \cup B$ é contado uma vez se estiver apenas em $A$ ou apenas em $B$. Porém, cada elemento de $A \cap B$ é contado \emph{duas vezes} --- uma vez ao contar $A$ e outra ao contar $B$. Para corrigir esta dupla contagem, subtraímos $n(A \cap B)$ uma vez:
\[
n(A \cup B) = n(A) + n(B) - n(A \cap B)
\]
\hfill$\square$

% --------------------------------------------------
\section*{Exercício 1.17}
\addcontentsline{toc}{subsection}{Exercício 1.17}
% --------------------------------------------------

\textbf{Enunciado:} Seja $\mathcal{A} = [\{1,2,3\}, \{4,5\}, \{6,7,8\}]$. (a) Listar elementos; (b) Encontrar $n(\mathcal{A})$.

\begin{teoribox}
\textbf{Classes de conjuntos:} Uma classe (ou família) de conjuntos é um conjunto cujos elementos são eles próprios conjuntos. A cardinalidade de uma classe conta o número de conjuntos que a compõem, não o número de elementos desses conjuntos.
\end{teoribox}

\textbf{(a)} $\mathcal{A}$ tem três elementos: os conjuntos $\{1,2,3\}$, $\{4,5\}$ e $\{6,7,8\}$.

\textbf{(b)} $n(\mathcal{A}) = 3$ (há três conjuntos na família, independentemente de quantos elementos eles têm internamente).

% --------------------------------------------------
\section*{Exercício 1.18}
\addcontentsline{toc}{subsection}{Exercício 1.18}
% --------------------------------------------------

\textbf{Enunciado:} Determinar o conjunto das partes $\mathcal{P}(A)$ de $A = \{a,b,c,d\}$.

\begin{teoribox}
\textbf{Conjunto das Partes:} $\mathcal{P}(A)$ (ou $2^A$) é o conjunto de todos os subconjuntos de $A$, incluindo $\emptyset$ e o próprio $A$. Se $n(A) = n$, então $n(\mathcal{P}(A)) = 2^n$.
\end{teoribox}

Como $n(A) = 4$, o conjunto das partes tem $2^4 = 16$ elementos. Listando por cardinalidade crescente:

\begin{align*}
\mathcal{P}(A) = \{&\emptyset, \\
&\{a\}, \{b\}, \{c\}, \{d\}, \\
&\{a,b\}, \{a,c\}, \{a,d\}, \{b,c\}, \{b,d\}, \{c,d\}, \\
&\{a,b,c\}, \{a,b,d\}, \{a,c,d\}, \{b,c,d\}, \\
&\{a,b,c,d\}\}
\end{align*}

% --------------------------------------------------
\section*{Exercício 1.19}
\addcontentsline{toc}{subsection}{Exercício 1.19}
% --------------------------------------------------

\textbf{Enunciado:} Seja $S = \{a,b,c,d,e,f,g\}$. Determinar quais são partições de $S$.

\begin{teoribox}
\textbf{Partição de um conjunto:} Uma colecção $\{A_i\}$ de subconjuntos não vazios de $S$ é uma partição se e só se:
\begin{enumerate}[label=(\roman*)]
  \item Os subconjuntos são dois a dois disjuntos: $A_i \cap A_j = \emptyset$ para $i \neq j$.
  \item A reunião cobre $S$: $\bigcup_i A_i = S$ (todo elemento de $S$ pertence a alguma célula).
\end{enumerate}
\end{teoribox}

\textbf{(a)} $P_1 = [\{a,c,e\}, \{b\}, \{d,g\}]$: o elemento $f \in S$ não pertence a nenhuma célula. \textbf{Não é partição.}

\textbf{(b)} $P_2 = [\{a,e,g\}, \{c,d\}, \{b,e,f\}]$: o elemento $e$ aparece nas células $\{a,e,g\}$ e $\{b,e,f\}$. \textbf{Não é partição.}

\textbf{(c)} $P_3 = [\{a,b,e,g\}, \{c\}, \{d,f\}]$: verificando --- $a,b,e,g$ na 1ª célula; $c$ na 2ª; $d,f$ na 3ª. Todos os 7 elementos de $S$ aparecem exactamente uma vez. \textbf{É uma partição.}

\textbf{(d)} $P_4 = [\{a,b,c,d,e,f,g\}]$: partição em uma única célula que é o próprio $S$. \textbf{É uma partição (trivial).}

% --------------------------------------------------
\section*{Exercício 1.20}
\addcontentsline{toc}{subsection}{Exercício 1.20}
% --------------------------------------------------

\textbf{Enunciado:} Encontrar todas as partições de $S = \{a,b,c,d\}$.

\textbf{Resolução:} Organizamos por número de células:

\textbf{1 célula:} $[\{a,b,c,d\}]$ --- 1 partição.

\textbf{2 células:} Uma célula de tamanho 1 e outra de tamanho 3, ou duas células de tamanho 2. Há $\binom{4}{1} = 4$ formas de escolher a célula de tamanho 1, e $\binom{4}{2}/2! \cdot 2! = 3$ formas de dividir em dois pares. Total: $4 + 3 = 7$ partições.

\textbf{3 células:} Duas células de tamanho 1 e uma de tamanho 2. Há $\binom{4}{2} = 6$ formas de escolher o par. Total: 6 partições.

\textbf{4 células:} $[\{a\},\{b\},\{c\},\{d\}]$ --- 1 partição.

Total geral: $1 + 7 + 6 + 1 = \mathbf{15}$ partições distintas. (Este número é o 4º número de Bell, $B_4 = 15$.)

% --------------------------------------------------
\section*{Exercício 1.21}
\addcontentsline{toc}{subsection}{Exercício 1.21}
% --------------------------------------------------

\textbf{Enunciado:} Para $A_n = \{n, 2n, 3n, \ldots\}$ (múltiplos de $n$), calcular: (a) $A_3 \cap A_5$; (b) $A_4 \cap A_6$; (c) $\bigcup_{i \in Q} A_i$ onde $Q$ é o conjunto dos primos.

\begin{teoribox}
\textbf{Relação entre múltiplos e mínimo múltiplo comum (m.m.c.):} Um número pertence a $A_m \cap A_n$ se e só se é múltiplo de ambos $m$ e $n$, ou seja, é múltiplo do m.m.c.$(m,n)$. Logo $A_m \cap A_n = A_{\text{m.m.c.}(m,n)}$.
\end{teoribox}

\textbf{(a)} $A_3 \cap A_5$: m.m.c.$(3,5) = 15$ (já que $\text{m.d.c.}(3,5)=1$). Logo $A_3 \cap A_5 = A_{15}$.

\textbf{(b)} $A_4 \cap A_6$: m.m.c.$(4,6) = 12$. Logo $A_4 \cap A_6 = A_{12}$.

\textbf{(c)} $\bigcup_{p \in Q} A_p$ é o conjunto de todos os inteiros positivos que são múltiplos de pelo menos um primo. Pelo Teorema Fundamental da Aritmética, todo inteiro $n > 1$ tem pelo menos um factor primo, logo $n \in \bigcup_{p \in Q} A_p$. O único inteiro positivo que não é múltiplo de qualquer primo é o $1$. Portanto:
\[
\bigcup_{p \in Q} A_p = \{2, 3, 4, \ldots\} = \mathbb{N} \setminus \{1\}
\]

% --------------------------------------------------
\section*{Exercício 1.22}
\addcontentsline{toc}{subsection}{Exercício 1.22}
% --------------------------------------------------

\textbf{Enunciado:} Seja $\{A_i \mid i \in I\}$ uma família indexada. Provar que para qualquer $i_0 \in I$:
\[
\bigcap_{i \in I} A_i \subseteq A_{i_0} \subseteq \bigcup_{i \in I} A_i
\]

\begin{teoribox}
\textbf{Intuição:} A intersecção de toda a família é o subconjunto mais ``restrito'' (só os elementos comuns a \emph{todos}), enquanto a reunião é o mais ``abrangente'' (tudo o que está em \emph{algum}). Qualquer membro individual da família está entre estes dois extremos.
\end{teoribox}

\textbf{Demonstração da inclusão à esquerda} ($\bigcap_{i \in I} A_i \subseteq A_{i_0}$):

Seja $x \in \bigcap_{i \in I} A_i$. Por definição da intersecção indexada, $x \in A_i$ para \emph{todo} $i \in I$. Em particular, para $i = i_0$, temos $x \in A_{i_0}$. Logo $\bigcap_{i \in I} A_i \subseteq A_{i_0}$.

\textbf{Demonstração da inclusão à direita} ($A_{i_0} \subseteq \bigcup_{i \in I} A_i$):

Seja $y \in A_{i_0}$. Como $i_0 \in I$, o conjunto $A_{i_0}$ é um dos conjuntos na reunião. Logo $y \in \bigcup_{i \in I} A_i$. Portanto $A_{i_0} \subseteq \bigcup_{i \in I} A_i$. \hfill$\square$

\newpage

% =============================================================
\section*{Capítulo 2 --- Relações}
\addcontentsline{toc}{section}{Capítulo 2 --- Relações}
% =============================================================

% --------------------------------------------------
\section*{Exercício 2.1}
\addcontentsline{toc}{subsection}{Exercício 2.1}
% --------------------------------------------------

\textbf{Enunciado:} Dados $A=\{1,2\}$, $B=\{x,y,z\}$, $C=\{3,4\}$. Encontrar $A \times B \times C$.

\begin{teoribox}
\textbf{Produto Cartesiano:} $A \times B \times C = \{(a,b,c) \mid a \in A,\, b \in B,\, c \in C\}$. A cardinalidade satisfaz $n(A \times B \times C) = n(A) \cdot n(B) \cdot n(C)$.
\end{teoribox}

Para obter sistematicamente todos os triplos, fixamos cada $a \in A$, depois cada $b \in B$, e por fim cada $c \in C$:
\[
A \times B \times C = \{(1,x,3),(1,x,4),(1,y,3),(1,y,4),(1,z,3),(1,z,4),(2,x,3),(2,x,4),(2,y,3),(2,y,4),(2,z,3),(2,z,4)\}
\]
A cardinalidade é $2 \times 3 \times 2 = 12$, confirmando os 12 elementos listados.

% --------------------------------------------------
\section*{Exercício 2.2}
\addcontentsline{toc}{subsection}{Exercício 2.2}
% --------------------------------------------------

\textbf{Enunciado:} Encontrar $x$ e $y$ dados $(2x, x+y) = (6, 2)$.

\begin{teoribox}
\textbf{Igualdade de pares ordenados:} $(a,b) = (c,d)$ se e só se $a = c$ \emph{e} $b = d$ (coordenada a coordenada).
\end{teoribox}

Igualando as coordenadas:
\begin{align*}
2x &= 6 \Rightarrow x = 3 \\
x + y &= 2 \Rightarrow 3 + y = 2 \Rightarrow y = -1
\end{align*}

Verificação: $(2 \cdot 3,\, 3 + (-1)) = (6, 2)$ \checkmark

% --------------------------------------------------
\section*{Exercício 2.3}
\addcontentsline{toc}{subsection}{Exercício 2.3}
% --------------------------------------------------

\textbf{Enunciado:} Encontrar o número de relações de $A = \{a,b,c\}$ para $B = \{1,2\}$.

\begin{teoribox}
\textbf{Relação binária:} Uma relação de $A$ para $B$ é um \emph{subconjunto} de $A \times B$. O número total de subconjuntos de um conjunto de $n$ elementos é $2^n$.
\end{teoribox}

O produto $A \times B$ tem $n(A) \cdot n(B) = 3 \times 2 = 6$ elementos. Cada subconjunto de $A \times B$ define uma relação. Logo, o número de relações é:
\[
2^{n(A \times B)} = 2^6 = 64
\]

% --------------------------------------------------
\section*{Exercício 2.4}
\addcontentsline{toc}{subsection}{Exercício 2.4}
% --------------------------------------------------

\textbf{Enunciado:} Para $A=\{1,2,3,4\}$, $B=\{x,y,z\}$ e $R = \{(1,y),(1,z),(3,y),(4,x),(4,z)\}$: (a) matriz; (b) grafo; (c) inversa $R^{-1}$; (d) domínio e imagem.

\begin{teoribox}
\textbf{Representações de uma relação:}
\begin{itemize}
  \item \textbf{Matriz:} Matriz booleana $M_R$ onde $(M_R)_{ij} = 1$ se $(a_i, b_j) \in R$, e $0$ caso contrário.
  \item \textbf{Relação inversa:} $R^{-1} = \{(b,a) \mid (a,b) \in R\}$ --- inverte cada par.
  \item \textbf{Domínio:} $\text{Dom}(R) = \{a \mid \exists b: (a,b) \in R\}$.
  \item \textbf{Imagem (Range):} $\text{Ran}(R) = \{b \mid \exists a: (a,b) \in R\}$.
\end{itemize}
\end{teoribox}

\textbf{(a) Matriz de $R$} (linhas = elementos de $A$, colunas = elementos de $B$):
\[
M_R = \begin{array}{c|ccc}
 & x & y & z \\ \hline
1 & 0 & 1 & 1 \\
2 & 0 & 0 & 0 \\
3 & 0 & 1 & 0 \\
4 & 1 & 0 & 1
\end{array}
\]

\textbf{(c) Inversa:} Inverte-se cada par ordenado de $R$:
\[
R^{-1} = \{(y,1),(z,1),(y,3),(x,4),(z,4)\}
\]

\textbf{(d) Domínio e imagem:}
\[
\text{Dom}(R) = \{1, 3, 4\}, \qquad \text{Ran}(R) = \{x, y, z\}
\]
Note que o elemento $2 \in A$ não aparece em nenhum par de $R$, pelo que $2 \notin \text{Dom}(R)$.

% --------------------------------------------------
\section*{Exercício 2.5}
\addcontentsline{toc}{subsection}{Exercício 2.5}
% --------------------------------------------------

\textbf{Enunciado:} Com $R = \{(1,b),(2,a),(2,c)\}$ e $S = \{(a,y),(b,x),(c,y),(c,z)\}$: (a) encontrar $R \circ S$; (b) comparar $M_{R \circ S}$ com $M_R M_S$.

\begin{teoribox}
\textbf{Composição de relações:} $(a,c) \in R \circ S$ se e só se existe $b$ tal que $(a,b) \in R$ e $(b,c) \in S$. O produto de matrizes booleanas $M_R M_S$ tem a mesma estrutura de zeros que $M_{R \circ S}$.
\end{teoribox}

\textbf{(a)} Para cada par $(a,b) \in R$, procuramos todos os $(b,c) \in S$:
\begin{align*}
(1,b) \in R &\text{ e } (b,x) \in S \Rightarrow (1,x) \in R \circ S \\
(2,a) \in R &\text{ e } (a,y) \in S \Rightarrow (2,y) \in R \circ S \\
(2,c) \in R &\text{ e } (c,y),(c,z) \in S \Rightarrow (2,y),(2,z) \in R \circ S
\end{align*}
Portanto: $R \circ S = \{(1,x),(2,y),(2,z)\}$.

\textbf{(b)} As matrizes são:
\[
M_R = \begin{pmatrix} 0&1&0 \\ 1&0&1 \\ 0&0&0 \end{pmatrix}, \quad
M_S = \begin{pmatrix} 0&1&0 \\ 1&0&0 \\ 0&1&1 \end{pmatrix}, \quad
M_{R \circ S} = \begin{pmatrix} 1&0&0 \\ 0&1&1 \\ 0&0&0 \end{pmatrix}
\]
O produto aritmético $M_R M_S$ produz:
\[
M_R M_S = \begin{pmatrix} 1&0&0 \\ 0&2&1 \\ 0&0&0 \end{pmatrix}
\]
Observa-se que $M_{R \circ S}$ e $M_R M_S$ têm exactamente as mesmas entradas \emph{nulas}, confirmando que os zeros do produto de matrizes correspondem às ausências de ligação na composição.

% --------------------------------------------------
\section*{Exercício 2.6}
\addcontentsline{toc}{subsection}{Exercício 2.6}
% --------------------------------------------------

\textbf{Enunciado:} Para $R = \{(1,1),(2,2),(2,3),(3,2),(4,2),(4,4)\}$ em $A = \{1,2,3,4\}$: (a) grafo dirigido; (b) $R^2 = R \circ R$.

\begin{teoribox}
\textbf{Composição $R^2$:} $(a,c) \in R^2$ se existe $b$ com $(a,b) \in R$ e $(b,c) \in R$.
\end{teoribox}

\textbf{(b)} Para cada $(a,b) \in R$, procuramos todos $(b,c) \in R$:
\begin{align*}
(1,1): &\quad (1,1) \in R \Rightarrow (1,1) \in R^2 \\
(2,2): &\quad (2,2),(2,3) \in R \Rightarrow (2,2),(2,3) \in R^2 \\
(2,3): &\quad (3,2) \in R \Rightarrow (2,2) \in R^2 \text{ (já contado)} \\
(3,2): &\quad (2,2),(2,3) \in R \Rightarrow (3,2),(3,3) \in R^2 \\
(4,2): &\quad (2,2),(2,3) \in R \Rightarrow (4,2),(4,3) \in R^2 \\
(4,4): &\quad (4,4) \in R \Rightarrow (4,4) \in R^2
\end{align*}
Portanto: $R^2 = \{(1,1),(2,2),(2,3),(3,2),(3,3),(4,2),(4,3),(4,4)\}$.

% --------------------------------------------------
\section*{Exercício 2.7}
\addcontentsline{toc}{subsection}{Exercício 2.7}
% --------------------------------------------------

\textbf{Enunciado:} Com $R = \{(1,1),(1,2),(2,3),(3,1),(3,3)\}$ e $S = \{(1,2),(1,3),(2,1),(3,3)\}$ em $A = \{1,2,3\}$. Calcular $R \cup S$, $R \cap S$, $R^c$, $R \circ S$, $S^2$.

\textbf{(a) Operações de conjunto em $R$ e $S$:}
\begin{align*}
R \cap S &= \{(1,2),(3,3)\} \\
R \cup S &= \{(1,1),(1,2),(1,3),(2,1),(2,3),(3,1),(3,3)\} \\
R^c &= (A \times A) \setminus R = \{(1,3),(2,1),(2,2),(3,2)\}
\end{align*}
Para $R^c$, o universo é $A \times A$ com 9 pares. $R$ tem 5 pares, logo $R^c$ tem $9 - 5 = 4$ pares.

\textbf{(b) Composição $R \circ S$:}

Para cada $(a,b) \in R$, encontramos $(b,c) \in S$:
\begin{align*}
(1,1) \in R, (1,2),(1,3) \in S &\Rightarrow (1,2),(1,3) \in R \circ S \\
(1,2) \in R, (2,1) \in S &\Rightarrow (1,1) \in R \circ S \\
(2,3) \in R, (3,3) \in S &\Rightarrow (2,3) \in R \circ S \\
(3,1) \in R, (1,2),(1,3) \in S &\Rightarrow (3,2),(3,3) \in R \circ S \\
(3,3) \in R, (3,3) \in S &\Rightarrow (3,3) \in R \circ S
\end{align*}
$R \circ S = \{(1,1),(1,2),(1,3),(2,3),(3,2),(3,3)\}$

\textbf{(c)} $S^2 = S \circ S = \{(1,1),(1,3),(2,2),(2,3),(3,3)\}$

% --------------------------------------------------
\section*{Exercício 2.8}
\addcontentsline{toc}{subsection}{Exercício 2.8}
% --------------------------------------------------

\textbf{Enunciado:} Provar o Teorema 2.1: $(R \circ S) \circ T = R \circ (S \circ T)$ (associatividade da composição).

\begin{teoribox}
\textbf{Estratégia:} Prova-se dupla inclusão. Para cada direcção, segue-se a cadeia de existência dos ``intermediários''.
\end{teoribox}

\textbf{Demonstração} ($\subseteq$): Suponha $(a,d) \in (R \circ S) \circ T$. Então existe $c$ tal que:
\[
(a,c) \in R \circ S \quad \text{e} \quad (c,d) \in T
\]
Como $(a,c) \in R \circ S$, existe $b$ tal que $(a,b) \in R$ e $(b,c) \in S$.

Como $(b,c) \in S$ e $(c,d) \in T$, temos $(b,d) \in S \circ T$.

Como $(a,b) \in R$ e $(b,d) \in S \circ T$, temos $(a,d) \in R \circ (S \circ T)$.

Logo $(R \circ S) \circ T \subseteq R \circ (S \circ T)$. A outra inclusão é simétrica. \hfill$\square$

% --------------------------------------------------
\section*{Exercício 2.9}
\addcontentsline{toc}{subsection}{Exercício 2.9}
% --------------------------------------------------

\textbf{Enunciado:} Para cinco relações em $A = \{1,2,3\}$, determinar reflexividade, simetria, transitividade e antisimetria.

\begin{teoribox}
\textbf{Propriedades de relações em $A$:}
\begin{itemize}
  \item \textbf{Reflexiva:} $(a,a) \in R$ para todo $a \in A$.
  \item \textbf{Simétrica:} $(a,b) \in R \Rightarrow (b,a) \in R$.
  \item \textbf{Transitiva:} $(a,b) \in R$ e $(b,c) \in R \Rightarrow (a,c) \in R$.
  \item \textbf{Antissimétrica:} $(a,b) \in R$ e $(b,a) \in R \Rightarrow a = b$.
\end{itemize}
\end{teoribox}

Relações: $R = \{(1,1),(1,2),(1,3),(3,3)\}$, $S = \{(1,1),(1,2),(2,1),(2,2),(3,3)\}$, $T = \{(1,1),(1,2),(2,2),(2,3)\}$, $\emptyset$, $A \times A$.

\textbf{(a) Reflexividade:} Precisamos de $(1,1),(2,2),(3,3)$ em cada relação.
\begin{itemize}
  \item $R$: falta $(2,2)$ --- \textbf{não reflexiva}.
  \item $S$: tem $(1,1),(2,2),(3,3)$ --- \textbf{reflexiva}.
  \item $T$: falta $(3,3)$ --- \textbf{não reflexiva}.
  \item $\emptyset$: faltam todos os pares diagonais --- \textbf{não reflexiva}.
  \item $A \times A$: contém todos os pares possíveis --- \textbf{reflexiva}.
\end{itemize}

\textbf{(b) Simetria:}
\begin{itemize}
  \item $R$: $(1,2) \in R$ mas $(2,1) \notin R$ --- \textbf{não simétrica}.
  \item $S$: para cada par $(a,b) \in S$ temos $(b,a) \in S$ --- \textbf{simétrica}.
  \item $T$: $(1,2) \in T$ mas $(2,1) \notin T$ --- \textbf{não simétrica}.
  \item $\emptyset$: vacuamente \textbf{simétrica} (não há pares para violar a condição).
  \item $A \times A$: contém todos os pares --- \textbf{simétrica}.
\end{itemize}

\textbf{(c) Transitividade:}
\begin{itemize}
  \item $R$: verificar todas as cadeias: $(1,2)$ e $(2,?)$ --- não há $(2,x) \in R$. $(1,3)$ e $(3,3) \Rightarrow (1,3) \in R$ \checkmark. \textbf{Transitiva}.
  \item $S$: \textbf{transitiva} (verificação detalhada: cada cadeia $a \to b \to c$ em $S$ tem $(a,c) \in S$).
  \item $T$: $(1,2),(2,3) \in T$ mas $(1,3) \notin T$ --- \textbf{não transitiva}.
  \item $\emptyset$: vacuamente \textbf{transitiva}.
  \item $A \times A$: \textbf{transitiva}.
\end{itemize}

\textbf{(d) Antisimetria:}
\begin{itemize}
  \item $S$: $1 \neq 2$, mas $(1,2),(2,1) \in S$ --- \textbf{não antissimétrica}.
  \item $A \times A$: igualmente --- \textbf{não antissimétrica}.
  \item $R$, $T$, $\emptyset$: \textbf{antissimétricas} (não existem pares $a \neq b$ com $(a,b)$ e $(b,a)$ ambos presentes).
\end{itemize}

% --------------------------------------------------
\section*{Exercício 2.10}
\addcontentsline{toc}{subsection}{Exercício 2.10}
% --------------------------------------------------

\textbf{Enunciado:} Dar exemplos de relações em $A = \{1,2,3\}$ com propriedades específicas.

\textbf{(a) Simétrica e antissimétrica:} $R = \{(1,1),(2,2)\}$.

Simétrica: $(1,1) \Rightarrow (1,1)$ \checkmark; $(2,2) \Rightarrow (2,2)$ \checkmark. Antissimétrica: os únicos pares têm $a = b$, logo a condição é satisfeita vacuamente para $a \neq b$.

\textbf{(b) Nem simétrica nem antissimétrica:} $R = \{(1,2),(2,3)\}$.

Não simétrica: $(1,2) \in R$ mas $(2,1) \notin R$. Não antissimétrica: não tem pares $(a,b)$ e $(b,a)$ com $a \neq b$ simultaneamente, mas também não é simétrica.

\begin{notabox}
\textit{Nota:} Simetria e antisimetria não são opostos! Uma relação pode ser ambas (ex.: a identidade) ou nenhuma.
\end{notabox}

\textbf{(c) Transitiva mas $R \cup R^{-1}$ não transitiva:} $R = \{(1,2)\}$.

$R$ é transitiva (não há cadeias de comprimento 2). $R^{-1} = \{(2,1)\}$. Então $R \cup R^{-1} = \{(1,2),(2,1)\}$, que \emph{não} é transitiva pois $(1,2),(2,1) \in R \cup R^{-1}$ mas $(1,1) \notin R \cup R^{-1}$.

% --------------------------------------------------
\section*{Exercício 2.11}
\addcontentsline{toc}{subsection}{Exercício 2.11}
% --------------------------------------------------

\textbf{Enunciado:} Seja $\mathcal{C}$ uma colecção de relações em $A$ e $T = \bigcap \{S \mid S \in \mathcal{C}\}$. Provar: (a) se todo $S$ é simétrico, $T$ é simétrico; (b) se todo $S$ é transitivo, $T$ é transitivo.

\textbf{(a) Simetria de $T$:}

Suponha $(a,b) \in T$. Por definição da intersecção, $(a,b) \in S$ para todo $S \in \mathcal{C}$. Como cada $S$ é simétrico, $(b,a) \in S$ para todo $S$. Logo $(b,a) \in \bigcap S = T$. \hfill$\square$

\textbf{(b) Transitividade de $T$:}

Suponha $(a,b) \in T$ e $(b,c) \in T$. Então $(a,b),(b,c) \in S$ para todo $S \in \mathcal{C}$. Como cada $S$ é transitivo, $(a,c) \in S$ para todo $S$. Logo $(a,c) \in T$. \hfill$\square$

% --------------------------------------------------
\section*{Exercício 2.12}
\addcontentsline{toc}{subsection}{Exercício 2.12}
% --------------------------------------------------

\textbf{Enunciado:} Seja $P$ uma propriedade de relações ``$R$-fechável'' se: (1) existe uma $P$-relação $S \supseteq R$; (2) a intersecção de $P$-relações é uma $P$-relação. (a) Mostrar que simetria e transitividade são $R$-fecháveis; (b) definir $P(R)$.

\textbf{(a)} A relação universal $A \times A$ é tanto simétrica como transitiva, e satisfaz $R \subseteq A \times A$ para qualquer $R$. Isso garante (1). O Exercício 2.11 garante (2). Logo, ambas as propriedades são $R$-fecháveis.

\textbf{(b)} Definindo $P(R) = \bigcap \{S \mid S \text{ é uma } P\text{-relação e } R \subseteq S\}$:

Este conjunto é não vazio (por (1)) e é uma $P$-relação (por (2)). Contém $R$ (pois cada $S$ na intersecção contém $R$). É a \emph{menor} $P$-relação contendo $R$ (a $P$-closure de $R$).

% --------------------------------------------------
\section*{Exercício 2.13}
\addcontentsline{toc}{subsection}{Exercício 2.13}
% --------------------------------------------------

\textbf{Enunciado:} Para $R = \{(a,a),(a,b),(b,c),(c,c)\}$ em $A = \{a,b,c\}$, encontrar os fechos reflexivo, simétrico e transitivo.

\begin{teoribox}
\textbf{Fechos de uma relação:}
\begin{itemize}
  \item \textbf{Fecho reflexivo:} $R \cup \{(a,a) \mid a \in A\}$ --- adicionar os pares diagonais em falta.
  \item \textbf{Fecho simétrico:} $R \cup R^{-1}$ --- adicionar o inverso de cada par.
  \item \textbf{Fecho transitivo:} $R \cup R^2 \cup R^3 \cup \cdots$ (para $|A| = n$, basta $R \cup R^2 \cup \cdots \cup R^n$).
\end{itemize}
\end{teoribox}

\textbf{(a) Fecho reflexivo:}

$R$ já contém $(a,a)$ e $(c,c)$ mas não $(b,b)$. Adicionando $(b,b)$:
\[
\text{reflexivo}(R) = \{(a,a),(a,b),(b,b),(b,c),(c,c)\}
\]

\textbf{(b) Fecho simétrico:}

$R^{-1} = \{(a,a),(b,a),(c,b),(c,c)\}$. Os pares novos são $(b,a)$ e $(c,b)$:
\[
\text{simétrico}(R) = \{(a,a),(a,b),(b,a),(b,c),(c,b),(c,c)\}
\]

\textbf{(c) Fecho transitivo:}

Calculamos $R^2 = R \circ R$: para cada $(x,y) \in R$, procuramos $(y,z) \in R$:
\begin{align*}
(a,a),(a,b) \in R &\Rightarrow (a,a),(a,b) \in R^2 \\
(a,b),(b,c) \in R &\Rightarrow (a,c) \in R^2 \\
(b,c),(c,c) \in R &\Rightarrow (b,c) \in R^2 \\
(c,c),(c,c) \in R &\Rightarrow (c,c) \in R^2
\end{align*}
$R^2 = \{(a,a),(a,b),(a,c),(b,c),(c,c)\}$

$R^3 = R^2 \circ R = \{(a,a),(a,b),(a,c),(b,c),(c,c)\}$ (igual a $R^2$, pois estabilizou).

\[
\text{transitivo}(R) = R \cup R^2 \cup R^3 = \{(a,a),(a,b),(a,c),(b,c),(c,c)\}
\]

% --------------------------------------------------
\section*{Exercício 2.14}
\addcontentsline{toc}{subsection}{Exercício 2.14}
% --------------------------------------------------

\textbf{Enunciado:} Mostrar que $x \equiv y \pmod{m}$ define uma relação de equivalência em $\mathbb{Z}$.

\begin{teoribox}
\textbf{Relação de Equivalência:} Uma relação $R$ é de equivalência se for \emph{reflexiva}, \emph{simétrica} e \emph{transitiva}. Definição: $x \equiv y \pmod{m}$ significa $m \mid (x - y)$, i.e., $x - y = km$ para algum $k \in \mathbb{Z}$.
\end{teoribox}

\textbf{(i) Reflexividade:} $x - x = 0 = 0 \cdot m$, logo $m \mid 0$. Assim $x \equiv x \pmod{m}$.

\textbf{(ii) Simetria:} Se $x \equiv y \pmod{m}$, então $m \mid (x - y)$, i.e., $x - y = km$. Logo $y - x = (-k)m$, pelo que $m \mid (y - x)$, i.e., $y \equiv x \pmod{m}$.

\textbf{(iii) Transitividade:} Se $x \equiv y \pmod{m}$ e $y \equiv z \pmod{m}$, então $x - y = km$ e $y - z = lm$. Somando:
\[
x - z = (x - y) + (y - z) = km + lm = (k + l)m
\]
Logo $m \mid (x - z)$, i.e., $x \equiv z \pmod{m}$. \hfill$\square$

% --------------------------------------------------
\section*{Exercício 2.15}
\addcontentsline{toc}{subsection}{Exercício 2.15}
% --------------------------------------------------

\textbf{Enunciado:} Em $A \times A$ (com $A$ conjunto de inteiros não nulos), definir $(a,b) \approx (c,d)$ se $ad = bc$. Provar que $\approx$ é de equivalência.

\begin{notabox}
\textit{Motivação:} Esta relação captura a igualdade de fracções: $(a,b)$ representa $\frac{a}{b}$, e $\frac{a}{b} = \frac{c}{d}$ se e só se $ad = bc$.
\end{notabox}

\textbf{(i) Reflexividade:} $(a,b) \approx (a,b)$ pois $ab = ba$ (comutatividade da multiplicação).

\textbf{(ii) Simetria:} Se $(a,b) \approx (c,d)$, então $ad = bc$. Por comutatividade, $cb = da$, logo $(c,d) \approx (a,b)$.

\textbf{(iii) Transitividade:} Se $(a,b) \approx (c,d)$ e $(c,d) \approx (e,f)$, então $ad = bc$ e $cf = de$. Multiplicando:
\[
(ad)(cf) = (bc)(de) \Rightarrow af \cdot cd = be \cdot cd
\]
Como $c \neq 0$ e $d \neq 0$, podemos cancelar $cd$:
\[
af = be \Rightarrow (a,b) \approx (e,f)
\]
\hfill$\square$

% --------------------------------------------------
\section*{Exercício 2.16}
\addcontentsline{toc}{subsection}{Exercício 2.16}
% --------------------------------------------------

\textbf{Enunciado:} Dada a relação de equivalência $R$ em $A = \{1,2,3,4,5,6\}$, encontrar a partição induzida.

\begin{teoribox}
\textbf{Classes de equivalência:} Para uma relação de equivalência $R$ em $A$, a classe de equivalência de $a$ é $[a] = \{x \in A \mid (a,x) \in R\}$. O conjunto quociente $A/R = \{[a] \mid a \in A\}$ forma uma partição de $A$.
\end{teoribox}

\textbf{Método:} Para encontrar $[a]$, olhamos para todos os $x$ tais que $(a,x) \in R$.

Da relação $R$ dada, verificamos:
\begin{itemize}
  \item Pares começando em $1$: $(1,1),(1,5)$ --- logo $[1] = \{1,5\}$.
  \item Pares começando em $2$: $(2,2),(2,3),(2,6)$ --- logo $[2] = \{2,3,6\}$.
  \item O único elemento restante é $4$: $(4,4)$ --- logo $[4] = \{4\}$.
\end{itemize}

Partição de $A$ induzida por $R$: $\{\{1,5\},\{2,3,6\},\{4\}\}$.

% --------------------------------------------------
\section*{Exercício 2.17}
\addcontentsline{toc}{subsection}{Exercício 2.17}
% --------------------------------------------------

\textbf{Enunciado:} Provar o Teorema 2.6: se $R$ é de equivalência em $A$, então $A/R$ é uma partição. Especificamente: (i) $a \in [a]$; (ii) $[a] = [b] \Leftrightarrow (a,b) \in R$; (iii) $[a] \neq [b] \Rightarrow [a] \cap [b] = \emptyset$.

\textbf{(i)} Como $R$ é reflexiva, $(a,a) \in R$, logo $a \in [a]$.

\textbf{(ii) ($\Rightarrow$):} Suponha $(a,b) \in R$. Para qualquer $x \in [b]$: $(b,x) \in R$. Por transitividade com $(a,b) \in R$: $(a,x) \in R$, logo $x \in [a]$. Assim $[b] \subseteq [a]$. Por simetria de $R$: $(b,a) \in R$, e analogamente $[a] \subseteq [b]$. Logo $[a] = [b]$.

\textbf{($\Leftarrow$):} Se $[a] = [b]$, por (i) $b \in [b] = [a]$, logo $(a,b) \in R$.

\textbf{(iii)} Prova por contrapositivo: se $[a] \cap [b] \neq \emptyset$, existe $x \in [a] \cap [b]$. Logo $(a,x) \in R$ e $(b,x) \in R$. Por simetria $(x,b) \in R$, e por transitividade com $(a,x) \in R$: $(a,b) \in R$. Por (ii), $[a] = [b]$. \hfill$\square$

% --------------------------------------------------
\section*{Exercício 2.18}
\addcontentsline{toc}{subsection}{Exercício 2.18}
% --------------------------------------------------

\textbf{Enunciado:} A inclusão de conjuntos $\subseteq$ é uma ordem parcial?

\begin{teoribox}
\textbf{Ordem Parcial:} Uma relação é uma ordem parcial (ou \emph{poset}) se for reflexiva, antissimétrica e transitiva.
\end{teoribox}

Para quaisquer conjuntos $A$, $B$, $C$:
\begin{itemize}
  \item \textbf{Reflexividade:} $A \subseteq A$ --- sempre verdade.
  \item \textbf{Antisimetria:} Se $A \subseteq B$ e $B \subseteq A$, então $A = B$ (por axioma da extensionalidade).
  \item \textbf{Transitividade:} Se $A \subseteq B$ e $B \subseteq C$, então $A \subseteq C$.
\end{itemize}
\textbf{Sim}, a inclusão de conjuntos é uma ordem parcial.

% --------------------------------------------------
\section*{Exercício 2.19}
\addcontentsline{toc}{subsection}{Exercício 2.19}
% --------------------------------------------------

\textbf{Enunciado:} Em $\mathbb{Z}$, definir $aRb$ por $b = a^r$ para algum inteiro positivo $r$. Mostrar que $R$ é ordem parcial.

\textbf{(a) Reflexividade:} $a = a^1$, logo $aRa$ (tome $r = 1$).

\textbf{(b) Antisimetria:} Suponha $aRb$ e $bRa$, com $b = a^r$ e $a = b^s$. Então:
\[
a = b^s = (a^r)^s = a^{rs}
\]
Daí $a^{rs-1} = 1$. Há três casos:
\begin{itemize}
  \item $rs = 1$: como $r,s$ são inteiros positivos, $r = s = 1$, logo $a = b$.
  \item $a = 1$: então $b = 1^r = 1 = a$.
  \item $a = -1$: então $b = (-1)^r \in \{-1,1\}$; como $b = a = -1$ ou $b = 1$ e $1 = (-1)^s$ implica $s$ par, mas então $a = 1^s = 1 \neq -1$, contradição. Logo $b = -1 = a$.
\end{itemize}
Em todos os casos $a = b$.

\textbf{(c) Transitividade:} Se $b = a^r$ e $c = b^s$, então $c = (a^r)^s = a^{rs}$. Como $rs$ é inteiro positivo, $aRc$. \hfill$\square$

\newpage

% =============================================================
\section*{Capítulo 3 --- Funções e Algoritmos}
\addcontentsline{toc}{section}{Capítulo 3 --- Funções e Algoritmos}
% =============================================================

% --------------------------------------------------
\section*{Exercício 3.1}
\addcontentsline{toc}{subsection}{Exercício 3.1}
% --------------------------------------------------

\textbf{Enunciado:} Seja $X = \{1,2,3,4\}$. Determinar se cada relação é uma função de $X$ em $X$.

\begin{teoribox}
\textbf{Definição de função:} Um subconjunto $f \subseteq X \times X$ é uma função $f: X \to X$ se e só se cada elemento $a \in X$ aparece como \emph{primeira coordenada em exactamente um} par ordenado em $f$. Isto é:
\begin{enumerate}[label=(\roman*)]
  \item \textbf{Existência (totalidade):} todo $a \in X$ aparece em algum par.
  \item \textbf{Unicidade:} não existem dois pares distintos com a mesma primeira coordenada.
\end{enumerate}
\end{teoribox}

\textbf{(a)} $f = \{(2,3),(1,4),(2,1),(3,2),(4,4)\}$

O elemento $2$ aparece em dois pares distintos: $(2,3)$ e $(2,1)$. Viola a unicidade. \textbf{Não é função.}

\textbf{(b)} $g = \{(3,1),(4,2),(1,1)\}$

O elemento $2 \in X$ não aparece como primeira coordenada em nenhum par de $g$. Viola a totalidade. \textbf{Não é função.}

\textbf{(c)} $h = \{(2,1),(3,4),(1,4),(2,1),(4,4)\}$

Embora $2$ apareça duas vezes como primeira coordenada, os pares são \emph{idênticos}: $(2,1) = (2,1)$. Portanto, como conjunto, $h = \{(2,1),(3,4),(1,4),(4,4)\}$. Cada elemento de $X$ aparece exactamente uma vez. \textbf{É função.}

% --------------------------------------------------
\section*{Exercício 3.2}
\addcontentsline{toc}{subsection}{Exercício 3.2}
% --------------------------------------------------

\textbf{Enunciado:} Esboçar o gráfico de: (a) $f(x) = x^2 + x - 6$;\quad (b) $g(x) = x^3 - 3x^2 - x + 3$.

\begin{teoribox}
\textbf{Método:} Para esboçar um gráfico, calculamos pares $(x, f(x))$ para vários valores de $x$, plotamos os pontos no plano cartesiano e traçamos uma curva suave. Para polinómios, é útil encontrar as raízes (zeros) e analisar o comportamento assimptótico.
\end{teoribox}

\textbf{(a) Raízes de $f(x) = x^2 + x - 6$:}

Factorizando: $f(x) = (x+3)(x-2)$. Logo as raízes são $x = -3$ e $x = 2$.

O vértice da parábola fica em $x = -\frac{1}{2}$ com $f(-\frac{1}{2}) = \frac{1}{4} - \frac{1}{2} - 6 = -\frac{25}{4}$.

Tabela de valores relevantes:
\[
\begin{array}{c|cccccc}
x & -4 & -3 & -1 & 0 & 2 & 3 \\\hline
f(x) & 6 & 0 & -6 & -6 & 0 & 6
\end{array}
\]

\textbf{(b) Raízes de $g(x) = x^3 - 3x^2 - x + 3$:}

Factorizando por agrupamento:
\[
g(x) = x^2(x-3) - 1(x-3) = (x^2-1)(x-3) = (x-1)(x+1)(x-3)
\]
Raízes: $x = -1$, $x = 1$, $x = 3$.

Tabela de valores:
\[
\begin{array}{c|cccccc}
x & -2 & -1 & 0 & 1 & 3 & 4 \\\hline
g(x) & -15 & 0 & 3 & 0 & 0 & 15
\end{array}
\]

% --------------------------------------------------
\section*{Exercício 3.3}
\addcontentsline{toc}{subsection}{Exercício 3.3}
% --------------------------------------------------

\textbf{Enunciado:} Com $f = \{(a,y),(b,x),(c,y)\}$ e $g = \{(x,s),(y,t),(z,r)\}$, encontrar (a) $g \circ f$; (b) $\text{Im}(f)$, $\text{Im}(g)$, $\text{Im}(g \circ f)$.

\begin{teoribox}
\textbf{Composição de funções:} $(g \circ f)(a) = g(f(a))$ --- aplica-se primeiro $f$, depois $g$ ao resultado. \textbf{Imagem (Im):} o conjunto de todos os valores efectivamente atingidos pela função.
\end{teoribox}

\textbf{(a)} Calculamos $(g \circ f)$ elemento a elemento:
\begin{align*}
(g \circ f)(a) &= g(f(a)) = g(y) = t \\
(g \circ f)(b) &= g(f(b)) = g(x) = s \\
(g \circ f)(c) &= g(f(c)) = g(y) = t
\end{align*}
Logo $g \circ f = \{(a,t),(b,s),(c,t)\}$.

\textbf{(b) Imagens:}
\begin{align*}
\text{Im}(f) &= \{y, x\} = \{x, y\} \quad \text{(apenas os valores atingidos por } f\text{; } z \text{ não é atingido)} \\
\text{Im}(g) &= \{s, t, r\} \quad \text{(} g \text{ é sobrejectiva em } \{r,s,t\}\text{)} \\
\text{Im}(g \circ f) &= \{t, s\} = \{s, t\} \quad \text{(} r \text{ não é atingido pois } z \notin \text{Im}(f)\text{)}
\end{align*}

\begin{notabox}
\textit{Observação:} $\text{Im}(g \circ f) \subseteq \text{Im}(g)$, e a igualdade só ocorre quando $\text{Im}(f)$ cobre todo o domínio de $g$.
\end{notabox}

% --------------------------------------------------
\section*{Exercício 3.4}
\addcontentsline{toc}{subsection}{Exercício 3.4}
% --------------------------------------------------

\textbf{Enunciado:} Com $f(x) = 2x+1$ e $g(x) = x^2-2$, encontrar a fórmula de $g \circ f$.

\begin{teoribox}
\textbf{Composição de funções reais:} $(g \circ f)(x) = g(f(x))$. Substitui-se a expressão de $f(x)$ no lugar de $x$ em $g$.
\end{teoribox}

\textbf{Método directo (substituição):}
\begin{align*}
(g \circ f)(x) &= g(f(x)) = g(2x+1) \\
&= (2x+1)^2 - 2 \\
&= 4x^2 + 4x + 1 - 2 \\
&= 4x^2 + 4x - 1
\end{align*}

\textbf{Método alternativo (eliminação de variável auxiliar):}

Defina $y = f(x) = 2x + 1$ e $z = g(y) = y^2 - 2$. Eliminando $y$:
\[
z = y^2 - 2 = (2x+1)^2 - 2 = 4x^2 + 4x - 1
\]

Ambos os métodos conduzem ao mesmo resultado:
\[
\boxed{(g \circ f)(x) = 4x^2 + 4x - 1}
\]

\begin{notabox}
\textit{Nota:} A composição de funções \textbf{não é comutativa} em geral. Comparar:
\[
(f \circ g)(x) = f(g(x)) = f(x^2 - 2) = 2(x^2 - 2) + 1 = 2x^2 - 3 \neq g \circ f
\]
\end{notabox}

% =============================================================
\newpage
\part*{Continuação --- Capítulos 10, 12, 13, 14, 15 e Apêndice B}
\addcontentsline{toc}{part}{Continuação --- Capítulos 10, 12, 13, 14, 15 e Apêndice B}
% =============================================================

%% ═════════════════════════════════════════════════════════════════════════════
\section*{Nota Introdutória}
\addcontentsline{toc}{section}{Nota Introdutória}

Este documento apresenta as resoluções detalhadas dos exercícios resolvidos dos capítulos
indicados, com ênfase didática: cada solução é desconstruída, os teoremas e propriedades
utilizados são identificados, e os cálculos intermédios são explicitados.
Os exercícios seguem a numeração original do livro.

\newpage
%% ═════════════════════════════════════════════════════════════════════════════
\part*{Capítulo 10 — Árvores Binárias}
\addcontentsline{toc}{part}{Capítulo 10 — Árvores Binárias}

%% ─────────────────────────────────────────────────────────────────────────────
\section*{Exercício 10.10 — Árvores Binárias Similares e Cópias}
\addcontentsline{toc}{section}{Exercício 10.10}

\begin{teoriacaixa}
\textbf{Conceitos essenciais.}
Duas árvores binárias $T$ e $T'$ dizem-se \keyword{similares} se têm a mesma estrutura
(a mesma ``forma''). Dizem-se \keyword{cópias} se são similares \emph{e} têm os mesmos
dados nos nós correspondentes. Numa árvore binária, distingue-se sempre o sucessor
\emph{esquerdo} do sucessor \emph{direito}, mesmo quando um nó tem apenas um filho.
\end{teoriacaixa}

\medskip
\textbf{Enunciado.} Considere as quatro árvores binárias da figura 10-1(b) do livro.
Indique quais são similares, quais são cópias, e justifique.

\medskip
\textbf{Resolução detalhada.}

As quatro árvores são rotuladas (1), (2), (3) e (4).

\begin{itemize}
  \item \textbf{Árvores (1), (3) e (4) são similares.}
  Para verificar a similaridade, compara-se recursivamente a estrutura: a raiz de cada uma
  tem dois filhos, que por sua vez têm (ou não) subárvores, de forma idêntica. A disposição
  topológica (quem é filho esquerdo, quem é filho direito) é a mesma nas três árvores.

  \item \textbf{Árvores (1) e (3) são cópias.}
  Além de serem similares, os dados armazenados nos nós correspondentes são idênticos.
  Formalmente, existe uma bijecção que preserva a estrutura \emph{e} os rótulos.

  \item \textbf{A árvore (2) \emph{não} é similar à árvore (4).}
  Embora (2) tenha o mesmo número de nós, numa das suas posições o único filho é
  o filho \emph{esquerdo}, enquanto na posição correspondente de (4) o único filho é
  o filho \emph{direito}. Numa árvore binária, esta distinção é fundamental:
  ``filho esquerdo'' e ``filho direito'' são entidades distintas, mesmo quando há
  apenas um filho. Por isso (2) e (4) diferem em estrutura, logo também não são cópias.
\end{itemize}

\medskip
\noindent\textit{Síntese:} A condição de similaridade exige igualdade estrutural completa,
incluindo a posição (esquerda/direita) de cada filho. A condição de cópia acrescenta a
igualdade de dados. A menor diferença de posição destrói a similaridade.

%% ─────────────────────────────────────────────────────────────────────────────
\section*{Exercício 10.11 — Representação Ligada de Árvores Binárias}
\addcontentsline{toc}{section}{Exercício 10.11}

\begin{teoriacaixa}
\textbf{Representação ligada (\textit{linked representation}).}
Uma árvore binária $T$ mantém-se em memória com três vetores paralelos \texttt{INFO},
\texttt{LEFT}, \texttt{RIGHT} e um ponteiro \texttt{ROOT}. Para cada nó $N$ que ocupa
a posição $K$:
\begin{align*}
  &\texttt{INFO}[K] \text{ --- dados do nó } N,\\
  &\texttt{LEFT}[K] \text{ --- localização do filho esquerdo de } N,\\
  &\texttt{RIGHT}[K] \text{ --- localização do filho direito de } N.
\end{align*}
Se uma subárvore é vazia, o ponteiro correspondente é \texttt{NULL} (convencionalmente
$0$ ou um número negativo).
\end{teoriacaixa}

\medskip
\textbf{Enunciado.} Dada a árvore $T$ da figura 10-1(a), a sua representação ligada
aparece na figura 10-5. Interpretá-la e justificar os valores.

\medskip
\textbf{Resolução detalhada.}

A árvore $T$ tem raiz $A$ e a representação satisfaz \texttt{ROOT} = 5, pois
$\texttt{INFO}[5] = A$: o nó raiz está armazenado na posição 5 dos vetores.

Para o nó $A$ (posição 5):
\[
  \texttt{LEFT}[5] = 10 \implies \texttt{INFO}[10] = B \quad \text{(filho esquerdo de } A \text{)},
\]
\[
  \texttt{RIGHT}[5] = 2 \implies \texttt{INFO}[2] = C \quad \text{(filho direito de } A \text{)}.
\]

Este esquema propaga-se recursivamente: para encontrar os filhos de qualquer nó $N$
na posição $K$, basta consultar $\texttt{LEFT}[K]$ e $\texttt{RIGHT}[K]$.

\medskip
\noindent\textbf{Porquê 18 elementos?}
A escolha de 18 posições para os vetores é arbitrária; o único requisito é que o número
de posições seja pelo menos igual ao número de nós da árvore. Posições não utilizadas
ficam vagas (ou marcadas como ``livres'').

%% ─────────────────────────────────────────────────────────────────────────────
\section*{Exercício 10.12 — Representação Sequencial de Árvores Binárias}
\addcontentsline{toc}{section}{Exercício 10.12}

\begin{teoriacaixa}
\textbf{Representação sequencial.}
Para uma árvore binária $T$ \emph{completa} (ou quase completa), é eficiente usar um
único vetor \texttt{TREE} com a convenção:
\begin{enumerate}[leftmargin=1.8em]
  \item A raiz $R$ fica em \texttt{TREE}[1].
  \item Se o nó $N$ ocupa \texttt{TREE}[$K$], então:\\
        filho esquerdo $\to$ \texttt{TREE}[$2K$],\\
        filho direito $\to$ \texttt{TREE}[$2K+1$].
  \item \texttt{END} guarda a posição do último nó.
\end{enumerate}
O pai do nó em \texttt{TREE}[$K$] está em \texttt{TREE}[$\lfloor K/2 \rfloor$].
\end{teoriacaixa}

\medskip
\textbf{Enunciado.} Para a árvore completa $T_{26}$ (figura 10-2), determine os filhos
do nó 9 e o pai do nó 9.

\medskip
\textbf{Resolução detalhada.}

Aplicando as fórmulas com $K = 9$:

\begin{align*}
  \text{Filho esquerdo} &= 2 \times 9 = 18,\\
  \text{Filho direito}  &= 2 \times 9 + 1 = 19,\\
  \text{Pai}            &= \lfloor 9/2 \rfloor = \lfloor 4{,}5 \rfloor = 4.
\end{align*}

\noindent\textbf{Verificação com a profundidade.} A profundidade $d_n$ da árvore
completa $T_n$ é:
\[
  d_n = \lfloor \log_2 n \rfloor + 1.
\]
Para $n = 26$: $\lfloor \log_2 26 \rfloor = \lfloor 4{,}70 \rfloor = 4$, logo $d_{26} = 5$.
Confirmamos que os nós 18 e 19 pertencem ao nível 4 (contando desde o nível 0 da raiz),
o que é consistente com a figura.

\newpage
%% ═════════════════════════════════════════════════════════════════════════════
\part*{Capítulo 12 — Linguagens, Autómatos e Gramáticas}
\addcontentsline{toc}{part}{Capítulo 12 — Linguagens, Autómatos e Gramáticas}

%% ─────────────────────────────────────────────────────────────────────────────
\section*{Exercícios 12.6--12.10 — Linguagens Formais}
\addcontentsline{toc}{section}{Exercícios 12.6--12.10}

\begin{teoriacaixa}
\textbf{Conceitos essenciais.}
Dado um alfabeto $A$, $A^*$ denota o conjunto de todas as palavras (sequências finitas)
sobre $A$, incluindo a \keyword{palavra vazia} $\lambda$.
Uma \keyword{linguagem} sobre $A$ é qualquer subconjunto $L \subseteq A^*$.
A \keyword{concatenação} de linguagens $K$ e $L$ é $KL = \{uv \mid u \in K,\, v \in L\}$.
A \keyword{potência} $L^n$ é a concatenação de $L$ consigo mesmo $n$ vezes ($L^0 = \{\lambda\}$).
O \keyword{fecho de Kleene} é $L^* = \bigcup_{n=0}^{\infty} L^n$.
\end{teoriacaixa}

\subsection*{Exercício 12.6}

\textbf{Enunciado.} Seja $A = \{a,b\}$. Descreva verbalmente:
$(a)$\; $L_1 = \{(ab)^m \mid m > 0\}$;\quad
$(b)$\; $L_2 = \{a^r b a^s b a^t \mid r,s,t \ge 0\}$;\quad
$(c)$\; $L_3 = \{a^2 b^m a^3 \mid m > 0\}$.

\medskip
\textbf{Resolução.}

\begin{enumerate}[label=(\alph*)]
  \item $L_1$ é o conjunto das palavras da forma $ab\,ab\,ab\cdots ab$, isto é, que
        começam com $a$, alternam entre $b$ e $a$ e terminam em $b$.
        Exemplos: $ab$, $abab$, $ababab$.

  \item $L_2$ contém exactamente todas as palavras com \emph{exactamente dois} $b$'s.
        Os expoentes $r,s,t \ge 0$ permitem qualquer número de $a$'s antes do primeiro
        $b$, entre os dois $b$'s, e depois do segundo $b$.
        Exemplo: $aaabab$ corresponde a $r=3,\,s=0,\,t=0$.

  \item $L_3$ são todas as palavras que começam com $a^2$, terminam com $a^3$, e têm
        pelo menos um $b$ no meio.
        Exemplo: $aabbba^3 = aabbbaaaa$.
\end{enumerate}

\subsection*{Exercício 12.7}

\textbf{Enunciado.} Sejam $K = \{a, ab, a^2\}$ e $L = \{b^2, aba\}$ linguagens sobre
$A = \{a,b\}$. Calcule $(a)$ $KL$ e $(b)$ $LL$.

\medskip
\textbf{Resolução.}

\begin{enumerate}[label=(\alph*)]
  \item Para obter $KL$, concatena-se cada elemento de $K$ com cada elemento de $L$:
  \begin{align*}
    a \cdot b^2 &= ab^2,\quad & a \cdot aba &= aaba,\\
    ab \cdot b^2 &= ab^3,\quad & ab \cdot aba &= ababa,\\
    a^2 \cdot b^2 &= a^2b^2,\quad & a^2 \cdot aba &= a^3ba.
  \end{align*}
  Portanto $KL = \{ab^2,\; a^2ba,\; ab^3,\; ababa,\; a^2b^2,\; a^3ba\}$.

  \item Para $LL$, concatenam-se os elementos de $L$ com os de $L$:
  \begin{align*}
    b^2 \cdot b^2 &= b^4,\quad & b^2 \cdot aba &= b^2aba,\\
    aba \cdot b^2 &= abab^2,\quad & aba \cdot aba &= aba^2ba.
  \end{align*}
  Portanto $LL = \{b^4,\; b^2aba,\; abab^2,\; aba^2ba\}$.
\end{enumerate}

\subsection*{Exercício 12.8}

\textbf{Enunciado.} Seja $L = \{ab, c\}$ sobre $A = \{a,b,c\}$. Calcule
$(a)$ $L^0$;\; $(b)$ $L^3$;\; $(c)$ $L^{-2}$.

\medskip
\textbf{Resolução.}

\begin{enumerate}[label=(\alph*)]
  \item Por definição, $L^0 = \{\lambda\}$, onde $\lambda$ é a palavra vazia.
        Esta convenção garante que $L^0 \cdot L = L$ para qualquer linguagem $L$.

  \item $L^3$ é o conjunto de todas as concatenações de \emph{três} palavras de $L$.
        Como $L$ tem dois elementos, existem $2^3 = 8$ combinações (mas verificar
        duplicados):
        \[
          L^3 = \{ababab,\; ababc,\; abcab,\; abc^2,\; cabab,\; cabc,\; c^2ab,\; c^3\}.
        \]
        Cada palavra resulta de escolher independentemente, para cada posição, $ab$ ou $c$.

  \item As potências negativas de uma linguagem \textbf{não estão definidas}.
        O conceito de ``inverso'' de uma concatenação não faz sentido no contexto das
        linguagens formais.
\end{enumerate}

\subsection*{Exercício 12.9}

\textbf{Enunciado.} Seja $A = \{a,b,c\}$. Descreva $L^*$ para:
$(a)$ $L = \{b^2\}$;\; $(b)$ $L = \{a,b\}$;\; $(c)$ $L = \{a,b,c^3\}$.

\medskip
\textbf{Resolução.}

O fecho de Kleene $L^*$ contém todas as palavras formadas por zero ou mais concatenações
de elementos de $L$.

\begin{enumerate}[label=(\alph*)]
  \item Cada elemento de $L^*$ é da forma $(b^2)^n = b^{2n}$ para $n \ge 0$.
        Logo $L^* = \{b^n \mid n \text{ é par}\} = \{\lambda, b^2, b^4, b^6, \ldots\}$.

  \item $L^*$ consiste em todas as palavras sobre $\{a,b\}$, pois $L = \{a,b\}$
        já é o alfabeto. Formalmente, $L^* = \{a,b\}^*$.

  \item As palavras em $L^*$ são concatenações de $a$'s, $b$'s e blocos $c^3$.
        A propriedade-chave é: qualquer subpalavra maximal composta exclusivamente de
        $c$'s tem comprimento divisível por 3. Isso porque cada ``contribuição'' de $c$'s
        vem do elemento $c^3 \in L$.
\end{enumerate}

\subsection*{Exercício 12.10}

\textbf{Enunciado.} Seja $A = \{a_1, a_2, \ldots\}$ um alfabeto contável. Seja $L_k$
a linguagem das palavras cuja soma dos índices é $k$.
$(a)$ Determine $L_4$.
$(b)$ Mostre que $L_k$ é finita.
$(c)$ Mostre que $A^*$ é contável.
$(d)$ Mostre que qualquer linguagem sobre $A$ é contável.

\medskip
\textbf{Resolução.}

\begin{enumerate}[label=(\alph*)]
  \item Para $k=4$: uma palavra em $L_4$ usa letras $a_i$ com $i \le 4$
        (pois um único $a_5$ já excede $k=4$), e tem no máximo 4 letras.
        Listando todas as formas de escrever 4 como soma ordenada de inteiros positivos:
        \[
          a_1a_1a_1a_1,\; a_1a_1a_2,\; a_1a_2a_1,\; a_2a_1a_1,\;
          a_1a_3,\; a_3a_1,\; a_2a_2,\; a_4.
        \]

  \item Para $k$ fixado, as letras usáveis são $a_1, \ldots, a_k$ (apenas $k$ letras),
        e cada palavra tem no máximo $k$ letras. Logo o número de palavras não excede
        $k^k$, que é finito. Portanto $L_k$ é finita.

  \item $A^* = \bigcup_{k=1}^{\infty} L_k$ é uma união contável de conjuntos finitos,
        logo contável. (União contável de conjuntos finitos é contável.)

  \item Qualquer linguagem $L \subseteq A^*$ é um subconjunto de um conjunto contável,
        logo também é contável (subconjuntos de conjuntos contáveis são contáveis).
\end{enumerate}

%% ─────────────────────────────────────────────────────────────────────────────
\section*{Exercícios 12.11--12.13 — Expressões Regulares}
\addcontentsline{toc}{section}{Exercícios 12.11--12.13}

\begin{teoriacaixa}
\textbf{Expressões regulares.} Uma \keyword{expressão regular} $r$ sobre $A$ é construída
recursivamente a partir de símbolos de $A$ usando os operadores:
\begin{itemize}[noitemsep]
  \item concatenação: $rs$ (justaposição),
  \item união (soma): $r \vee s$ (ou $r + s$),
  \item fecho de Kleene: $r^*$.
\end{itemize}
A linguagem $L(r)$ denotada por $r$ é definida pelas regras habituais.
\end{teoriacaixa}

\subsection*{Exercício 12.11}

\textbf{Enunciado.} Seja $A = \{a,b\}$. Descreva $L(r)$ para:
$(a)$ $r = abb^*a$;\;
$(b)$ $r = b^*ab^*ab^*$;\;
$(c)$ $r = a^* \vee b^*$;\;
$(d)$ $r = ab^* \wedge a^*$.

\medskip
\textbf{Resolução.}

\begin{enumerate}[label=(\alph*)]
  \item $r = abb^*a$: começa com $a$, seguido de $b$ (pelo menos um), depois $b^*$
        (zero ou mais $b$'s adicionais), terminando em $a$. Logo $L(r)$ contém todas as
        palavras que começam e terminam com $a$ e têm \emph{um ou mais} $b$'s no meio.

  \item $r = b^*ab^*ab^*$: a estrutura $b^*ab^*ab^*$ coloca exactamente dois $a$'s,
        com zero ou mais $b$'s em cada intervalo. Logo $L(r)$ são todas as palavras com
        \emph{exactamente dois} $a$'s.

  \item $r = a^* \vee b^*$: a união de $a^*$ e $b^*$ dá todas as palavras formadas
        exclusivamente por $a$'s \emph{ou} exclusivamente por $b$'s (incluindo $\lambda$).
        Assim $L(r) = \{\lambda, a, a^2, \ldots\} \cup \{\lambda, b, b^2, \ldots\}$.

  \item A expressão $ab^* \wedge a^*$ \textbf{não} é uma expressão regular, pois o
        símbolo $\wedge$ (conjunção) não faz parte dos operadores das expressões regulares.
        A classe das linguagens regulares é fechada para interseção, mas $\wedge$ não é
        um operador da notação standard de expressões regulares.
\end{enumerate}

\subsection*{Exercício 12.12}

\textbf{Enunciado.} Seja $A = \{a,b,c\}$ e $w = abc$. Determine se $w \in L(r)$ para:
$(a)$ $r = a^* \vee (b \vee c)^*$;\; $(b)$ $r = a^*(b \vee c)^*$.

\medskip
\textbf{Resolução.}

\begin{enumerate}[label=(\alph*)]
  \item $L(r) = L(a^*) \cup L((b \vee c)^*)$ contém palavras formadas \emph{só} por $a$'s,
        ou \emph{só} por $b$'s e $c$'s. A palavra $w = abc$ mistura $a$ com $b,c$,
        pelo que \textbf{não pertence} a $L(r)$.

  \item $r = a^*(b \vee c)^*$: primeiro um bloco (possivelmente vazio) de $a$'s,
        depois um bloco (possivelmente vazio) de $b$'s e $c$'s.
        Para $w = abc$: $a \in L(a^*)$ e $bc \in L((b \vee c)^*)$ (pois $b$ e $c$ são
        ambos produzidos por $b \vee c$). Logo $w = a \cdot bc \in L(r)$.
        Resposta: \textbf{sim}.
\end{enumerate}

\subsection*{Exercício 12.13}

\textbf{Enunciado.} Seja $A = \{a,b\}$. Encontre uma expressão regular $r$ tal que $L(r)$
consista em:
$(a)$ todas as palavras que começam com $a^2$ e terminam com $b^2$;
$(b)$ todas as palavras com um número par de $a$'s.

\medskip
\textbf{Resolução.}

\begin{enumerate}[label=(\alph*)]
  \item A expressão $r = a^2 (a \vee b)^* b^2$ funciona: $a^2$ força o início com dois
        $a$'s; $(a \vee b)^*$ aceita qualquer coisa (incluindo $\lambda$) no meio;
        $b^2$ força o fim com dois $b$'s. Logo
        \[
          r = a^2(a \vee b)^*b^2.
        \]
        Nota: são necessárias pelo menos 4 letras, pois $a^2b^2$ é a palavra mais curta.

  \item Palavras com número par de $a$'s incluem $0, 2, 4, \ldots$ $a$'s.
        A expressão $s = b^*ab^*ab^*$ captura exactamente as palavras com \emph{dois}
        $a$'s. O fecho $s^*$ gera concatenações de zero ou mais palavras com dois $a$'s
        cada, logo qualquer número \emph{par} de $a$'s, com $b$'s arbitrários. Portanto:
        \[
          r = s^* = (b^*ab^*ab^*)^*.
        \]
\end{enumerate}

%% ─────────────────────────────────────────────────────────────────────────────
\section*{Exercícios 12.14--12.21 — Autómatos Finitos}
\addcontentsline{toc}{section}{Exercícios 12.14--12.21}

\begin{teoriacaixa}
\textbf{Autómato finito determinístico (AFD).}
Um AFD é um quíntuplo $M = (A, S, Y, s_0, F)$ onde:
\begin{itemize}[noitemsep]
  \item $A$: alfabeto de entrada,
  \item $S$: conjunto finito de estados,
  \item $Y \subseteq S$: estados de aceitação,
  \item $s_0 \in S$: estado inicial,
  \item $F: S \times A \to S$: função de transição.
\end{itemize}
Uma palavra $w$ é \keyword{aceite} por $M$ se, ao processar $w$ a partir de $s_0$,
o autómato termina num estado de $Y$.
\end{teoriacaixa}

\subsection*{Exercício 12.14}

\textbf{Enunciado.} Seja $M$ com $A = \{a,b\}$, $S = \{s_0, s_1, s_2\}$, $Y = \{s_1\}$,
e a função de transição dada por:
\[
  F(s_0, a) = s_0,\quad F(s_0, b) = s_1,\quad
  F(s_1, a) = s_1,\quad F(s_1, b) = s_2,\quad
  F(s_2, a) = s_2,\quad F(s_2, b) = s_2.
\]
$(a)$ Desenhe o diagrama de estados. $(b)$ Descreva $L(M)$.

\medskip
\textbf{Resolução.}

\begin{enumerate}[label=(\alph*)]
  \item \textbf{Diagrama de estados:}
  Os vértices são os estados $s_0, s_1, s_2$ (com $s_1$ a círculo duplo por ser estado
  de aceitação). As transições são as arestas dirigidas rotuladas por $a$ ou $b$.
  Uma seta especial indica o estado inicial $s_0$.

  \item \textbf{Linguagem aceite.} Analisando a função de transição:
  \begin{itemize}
    \item $s_0$ é o estado ``nenhum $b$ visto ainda'': lendo $a$'s permanece em $s_0$;
          ao ler o primeiro $b$ vai para $s_1$.
    \item $s_1$ é o estado ``exactamente um $b$ lido'': $a$'s adicionais mantêm o estado;
          um segundo $b$ envia para $s_2$.
    \item $s_2$ é um estado ``sumidouro'' (dois ou mais $b$'s): nunca sai de $s_2$.
  \end{itemize}
  Logo $L(M)$ consiste em todas as palavras de $A^*$ com \textbf{exactamente um} $b$.
\end{enumerate}

\subsection*{Exercício 12.15}

\textbf{Enunciado.} Construa um autómato $M$ sobre $A = \{a,b\}$ que aceite precisamente
as palavras com um número par de $a$'s.

\medskip
\textbf{Resolução.}

A ideia é manter na memória (através dos estados) a \emph{paridade} do número de $a$'s
lidos. São necessários dois estados:

\begin{itemize}
  \item $s_0$: número par de $a$'s até ao momento (estado inicial e de aceitação).
  \item $s_1$: número ímpar de $a$'s até ao momento.
\end{itemize}

As transições são:
\[
  F(s_0, a) = s_1,\quad F(s_0, b) = s_0,\quad
  F(s_1, a) = s_0,\quad F(s_1, b) = s_1.
\]
O input $b$ não altera a paridade, logo não muda o estado. O input $a$ troca entre os
dois estados. O conjunto de aceitação é $Y = \{s_0\}$.

\medskip
\noindent\textit{Verificação:} a palavra $aababbab$ tem $a$'s nas posições 1, 2, 4, 7 —
quatro $a$'s (par) — pelo que deve ser aceite. Seguindo as transições:
$s_0 \xrightarrow{a} s_1 \xrightarrow{a} s_0 \xrightarrow{b} s_0 \xrightarrow{a} s_1
\xrightarrow{b} s_1 \xrightarrow{b} s_1 \xrightarrow{a} s_0 \xrightarrow{b} s_0 \in Y$. (aplicável)

\subsection*{Exercício 12.16}

\textbf{Enunciado.} Construa um autómato $M$ sobre $A = \{a,b\}$ que aceite palavras que
começam com $a$ seguido de zero ou mais $b$'s.

\medskip
\textbf{Resolução.}

São necessários três estados:

\begin{itemize}
  \item $s_0$: estado inicial (ainda não foi lida nenhuma letra).
  \item $s_1$: foi lido o $a$ inicial; lendo $b$'s adicionais permanece aqui (estado de aceitação).
  \item $s_2$: estado de rejeição (entrada inválida).
\end{itemize}

Transições:
\[
  F(s_0, a) = s_1,\quad F(s_0, b) = s_2,\quad
  F(s_1, a) = s_2,\quad F(s_1, b) = s_1,\quad
  F(s_2, a) = s_2,\quad F(s_2, b) = s_2.
\]
Conjunto de aceitação: $Y = \{s_1\}$.

A linguagem aceite é $L(M) = \{ab^n \mid n \ge 0\} = \{a, ab, abb, abbb, \ldots\}$.

\subsection*{Exercício 12.17}

\textbf{Enunciado.} Descreva as palavras $w$ aceites pelo autómato $M$ na figura 12-9(a).

\medskip
\textbf{Resolução.}

O autómato tem estados $\{s_0, s_1, s_2\}$ com $s_2$ como estado de aceitação.
Analisando as transições: o sistema só pode atingir $s_2$ quando lê um $a$ \emph{após}
ter lido um $b$. Mais precisamente:
\begin{itemize}
  \item $s_0$ representa ``nenhum $b$ foi lido ainda''; $a$ mantém em $s_0$; $b$ vai para $s_1$.
  \item $s_1$ representa ``foi lido pelo menos um $b$''; $a$ vai para $s_2$ (aceitação!);
        $b$ mantém em $s_1$.
  \item $s_2$ é estado de aceitação; $a$ e $b$ mantêm em $s_2$.
\end{itemize}

Logo $L(M)$ consiste em todas as palavras $w$ que contêm pelo menos um $a$ que \emph{sucede}
algum $b$, ou seja, $w$ contém a subpalavra $ba$ em algum ponto.

\subsection*{Exercício 12.18}

\textbf{Enunciado.} Descreva as palavras $w$ aceites pelo autómato $M$ na figura 12-9(b).

\medskip
\textbf{Resolução.}

O autómato tem quatro estados $\{s_0, s_1, s_2, s_3\}$ com $s_3$ como estado de aceitação.
As transições verificam-se:
\begin{itemize}
  \item Cada $a$ em $w$ \emph{não altera} o estado (laço em cada estado).
  \item Cada $b$ avança o estado: $s_0 \to s_1 \to s_2 \to s_3 \to s_0$ (cíclico módulo 4).
\end{itemize}

O estado atingido após ler $w$ depende apenas do \emph{número de $b$'s} em $w$,
digamos $n$. O estado final é $s_j$ onde $j \equiv n \pmod{4}$.
Como $Y = \{s_3\}$, a palavra é aceite se e só se $n \equiv 3 \pmod{4}$,
isto é, $n \in \{3, 7, 11, 15, \ldots\}$.

\subsection*{Exercício 12.19}

\textbf{Enunciado.} Dado $M = (A, S, Y, s_0, F)$ que aceita $L$, construa um autómato
$N$ que aceite $L^C$ (o complementar de $L$).

\medskip
\textbf{Resolução.}

A ideia é simples: uma palavra $w \notin L$ é exactamente uma palavra que $M$ \emph{rejeita},
ou seja, que termina num estado \emph{fora} de $Y$. Logo basta trocar os estados de
aceitação e rejeição.

Formalmente:
\[
  N = (A,\; S,\; S \setminus Y,\; s_0,\; F).
\]
$N$ usa o mesmo alfabeto, os mesmos estados, a mesma função de transição e o mesmo estado
inicial que $M$; apenas o conjunto de aceitação muda de $Y$ para $S \setminus Y$.

\medskip
\noindent\textit{Porquê funciona?} Para qualquer $w \in A^*$, a computação de $N$ em
$w$ é idêntica à de $M$ (mesma sequência de estados). Mas $N$ aceita exactamente quando
$M$ rejeita, e rejeita exactamente quando $M$ aceita.

\subsection*{Exercício 12.20}

\textbf{Enunciado.} Dados $M = (A, S, Y, s_0, F)$ e $M' = (A, S', Y', s_0', F')$,
construa $N$ que aceite $L(M) \cap L(M')$.

\medskip
\textbf{Resolução.}

Usa-se a \keyword{construção produto}: o novo autómato $N$ simula $M$ e $M'$ em
\emph{paralelo}, mantendo o par de estados como único estado.

\begin{itemize}
  \item Conjunto de estados de $N$: $S \times S'$.
  \item Estado inicial: $(s_0, s_0')$.
  \item Estados de aceitação: $Y \times Y' = \{(s, s') \mid s \in Y,\, s' \in Y'\}$;
        aceitar quando \emph{ambos} os autómatos aceitariam.
  \item Função de transição:
        $G\bigl((s, s'),\, a\bigr) = \bigl(F(s,a),\, F'(s', a)\bigr)$.
\end{itemize}

\medskip
\noindent\textit{Correctude:} uma palavra $w$ leva $N$ de $(s_0, s_0')$ para
$(F^*(s_0, w),\, F'^*(s_0', w))$. Este estado pertence a $Y \times Y'$ se e só se
$w \in L(M)$ \emph{e} $w \in L(M')$.

\subsection*{Exercício 12.21}

\textbf{Enunciado.} Dados $M$ e $M'$ como acima, construa $N$ que aceite $L(M) \cup L(M')$.

\medskip
\textbf{Resolução.}

A construção é idêntica à do Exercício 12.20, mas o conjunto de aceitação muda:

\[
  Y_N = \{(s, s') \in S \times S' \mid s \in Y \text{ ou } s' \in Y'\}.
\]

Ou seja, $N$ aceita quando \emph{pelo menos um} dos autómatos originais aceitaria.
Todo o resto (estados, estado inicial, função de transição) é idêntico à construção
produto do exercício anterior.

\newpage
%% ═════════════════════════════════════════════════════════════════════════════
\part*{Capítulo 13 — Máquinas de Estados Finitos e Máquinas de Turing}
\addcontentsline{toc}{part}{Capítulo 13 — Máquinas de Turing}

\begin{teoriacaixa}
\textbf{Máquina de Turing (MT).}
Uma MT é definida por quintuples $q = s_i a_j s_k a_l D$ onde:
$s_i$ é o estado actual, $a_j$ é o símbolo lido, $s_k$ é o novo estado,
$a_l$ é o símbolo escrito, e $D \in \{L, R, N\}$ é a direcção de movimento
(esquerda, direita, imóvel). Um \keyword{quadro} (picture) $\alpha$ descreve
a configuração instantânea: a fita, com o estado inserido na posição do cabeçote.
\end{teoriacaixa}

\section*{Exercício 13.1 (= 13.14 do livro) — Quadros de Máquinas de Turing}
\addcontentsline{toc}{section}{Exercício 13.1 (13.14)}

\textbf{Enunciado.} Determine o quadro $\alpha$ para cada situação:
$(a)$ A MT está no estado $s_2$ e lê o 3.º símbolo de $w = abbaa$.
$(b)$ A MT está no estado $s_3$ e lê o último símbolo de $w = aabb$.
$(c)$ O input é $W = a^3b^3$.
$(d)$ O input é $W = {)}3{,}2{)}*$.

\medskip
\textbf{Resolução.}

O quadro obtém-se inserindo o rótulo do estado imediatamente à \emph{esquerda} do símbolo
que está a ser lido.

\begin{enumerate}[label=(\alph*)]
  \item $w = abbaa$; 3.º símbolo é o segundo $b$ (posição 3). Inserindo $s_2$ antes desse $b$:
        \[
          \alpha = ab\;s_2\;baa.
        \]

  \item $w = aabb$; último símbolo é $b$ (posição 4). Inserindo $s_3$ antes:
        \[
          \alpha = aab\;s_3\;b.
        \]

  \item O input $W = a^3b^3 = aaabbb$. No início, a MT está no estado $s_0$ a ler o
        primeiro símbolo:
        \[
          \alpha = s_0\;aaabbb.
        \]

  \item A notação $)(3,2)*$ representa o par $(3,2)$ em código unário: $)3* = 1111$
        e $)2* = 111$, separados por $B$. Logo $)(3,2)* = 1111B111$.
        O quadro inicial é:
        \[
          \alpha = s_0\;1111B111.
        \]
\end{enumerate}

\section*{Exercício 13.2 (= 13.15 do livro) — Transições de Quadros}
\addcontentsline{toc}{section}{Exercício 13.2 (13.15)}

\textbf{Enunciado.} Seja $\alpha = abs_2aa$. Determine $\beta$ tal que $\alpha \to \beta$
para cada quintuple $q$:
$(a)$ $q = s_2 a b s_1 R$;\;
$(b)$ $q = s_2 a a s_3 L$;\;
$(c)$ $q = s_2 a b s_2 R$;\;
$(d)$ $q = s_2 a b s_3 L$;\;
$(e)$ $q = s_3 a b s_2 R$;\;
$(f)$ $q = s_2 a a s_2 L$.

\medskip
\textbf{Resolução.}

No quadro $\alpha = abs_2aa$, a MT está no estado $s_2$ a ler $a$ (o símbolo imediatamente
à direita de $s_2$). O quintuple $s_i a_j s_k a_l D$ só se aplica se $s_i$ e $a_j$
coincidirem com o estado e símbolo actuais.

\begin{enumerate}[label=(\alph*)]
  \item $q = s_2 a b s_1 R$: estado $s_2$, lê $a$ (aplicável). Substitui $a$ por $b$, vai para $s_1$,
        move à direita. O cabeçote fica sobre o próximo símbolo ($a$):
        \[
          \beta = ab\,b\,s_1\,a.
        \]

  \item $q = s_2 a a s_3 L$: estado $s_2$, lê $a$ (aplicável). Escreve $a$ (sem mudança), vai para $s_3$,
        move à esquerda. O cabeçote recua sobre $b$:
        \[
          \beta = a\,s_3\,b\,aa.
        \]

  \item $q = s_2 a b s_2 R$: estado $s_2$, lê $a$ (aplicável). Substitui $a$ por $b$, permanece em $s_2$,
        move à direita:
        \[
          \beta = ab\,b\,s_2\,a.
        \]

  \item $q = s_2 a b s_3 L$: estado $s_2$, lê $a$ (aplicável). Substitui $a$ por $b$, vai para $s_3$,
        move à esquerda (o cabeçote recua sobre $b$):
        \[
          \beta = a\,s_3\,b\,ba.
        \]

  \item $q = s_3 a b s_2 R$: estado actual é $s_2 \ne s_3$. O quintuple \emph{não se aplica};
        $\alpha$ não é alterado.

  \item $q = s_2 a a s_2 L$: estado $s_2$, lê $a$ (aplicável). Escreve $a$ (sem mudança), permanece em
        $s_2$, move à esquerda:
        \[
          \beta = a\,s_2\,b\,aa.
        \]
\end{enumerate}

\section*{Exercício 13.3 (= 13.17 do livro) — Ciclos Infinitos}
\addcontentsline{toc}{section}{Exercício 13.3 (13.17)}

\textbf{Enunciado.} Encontre quadros distintos $\alpha_1, \alpha_2, \alpha_3, \alpha_4$
e uma MT $M$ tais que a sequência não termine:
$\alpha_1 \to \alpha_2 \to \alpha_3 \to \alpha_4 \to \alpha_1 \to \cdots$

\medskip
\textbf{Resolução.}

Definem-se quatro quadros num ciclo de comprimento 4:

\begin{align*}
  \alpha_1 &= s_0 ab,\\
  \alpha_2 &= b\,s_1\,b,\\
  \alpha_3 &= s_2\,bb,\\
  \alpha_4 &= a\,s_3\,b.
\end{align*}

E os quintuples que produzem as transições:
\begin{align*}
  q_1 &= s_0\,a\,b\,s_1\,R &\text{($s_0$ lê $a$, escreve $b$, vai } s_1, R\text{)},\\
  q_2 &= s_1\,b\,b\,s_2\,L &\text{($s_1$ lê $b$, escreve $b$, vai } s_2, L\text{)},\\
  q_3 &= s_2\,b\,a\,s_3\,R &\text{($s_2$ lê $b$, escreve $a$, vai } s_3, R\text{)},\\
  q_4 &= s_3\,b\,b\,s_0\,L &\text{($s_3$ lê $b$, escreve $b$, vai } s_0, L\text{)}.
\end{align*}

Verificação: $\alpha_1 = s_0 ab \xrightarrow{q_1} bs_1b = \alpha_2 \xrightarrow{q_2}
s_2 bb = \alpha_3 \xrightarrow{q_3} as_3 b = \alpha_4 \xrightarrow{q_4} s_0 ab = \alpha_1$.
O ciclo repete indefinidamente.

\section*{Exercício 13.4 (= 13.18/13.19 do livro) — Determinismo e Reconhecimento}
\addcontentsline{toc}{section}{Exercício 13.4 (13.18--13.19)}

\textbf{Exercício 13.18.} Se $\alpha \to \beta_1$ e $\alpha \to \beta_2$, será que $\beta_1 = \beta_2$?

\medskip
\textbf{Resolução.} \textbf{Sim.} Uma Máquina de Turing é \keyword{determinística}: dado
o estado actual e o símbolo lido, existe no máximo um quintuple aplicável. Logo a próxima
configuração é \emph{única}. Se $\alpha \to \beta_1$ e $\alpha \to \beta_2$, então
necessariamente $\beta_1 = \beta_2$.

\medskip

\textbf{Exercício 13.19.} Se $\alpha = \alpha(W)$ para algum input $W$, e $\alpha \to \beta \to \alpha$,
pode $M$ reconhecer $W$?

\medskip
\textbf{Resolução.} \textbf{Não.} O ciclo $\alpha \to \beta \to \alpha \to \beta \to \cdots$
é infinito: a computação nunca termina. Uma MT reconhece $W$ apenas se a computação
\emph{pára} num estado $s_Y$ de aceitação. Como a computação entra num ciclo e nunca pára,
$W$ não é reconhecido.

\section*{Exercício 13.5 (= 13.20 do livro) — MT que Reconhece $L = \{ab^n \mid n > 0\}$}
\addcontentsline{toc}{section}{Exercício 13.5 (13.20)}

\textbf{Enunciado.} Encontre uma MT $M$ que reconhece a linguagem $L = \{ab^n \mid n > 0\}$,
ou seja, o conjunto das palavras que começam por $a$ seguido de um ou mais $b$'s.

\medskip
\textbf{Resolução.}

A estratégia é ler letra a letra e verificar:
\begin{enumerate}
  \item O primeiro símbolo tem de ser $a$.
  \item Depois deve haver pelo menos um $b$.
  \item Só devem existir $b$'s após o $a$ inicial.
  \item A palavra termina num branco $B$.
\end{enumerate}

Os quintuples (com $s_N$ = ``rejeita'' e $s_Y$ = ``aceita''):

\begin{align*}
  q_1 &= s_0\,B\,B\,s_N\,R &\text{(vazio: rejeitar)},\\
  q_2 &= s_0\,b\,b\,s_N\,R &\text{(começa com } b \text{: rejeitar)},\\
  q_3 &= s_0\,a\,a\,s_1\,R &\text{(leu } a \text{ inicial: avançar)},\\
  q_4 &= s_1\,B\,B\,s_N\,R &\text{(nenhum } b \text{ após } a \text{: rejeitar)},\\
  q_5 &= s_1\,a\,a\,s_N\,R &\text{(segundo } a \text{: rejeitar)},\\
  q_6 &= s_1\,b\,b\,s_2\,R &\text{(primeiro } b \text{: avançar para } s_2\text{)},\\
  q_7 &= s_2\,b\,b\,s_2\,R &\text{(continua a ler } b \text{'s)},\\
  q_8 &= s_2\,a\,a\,s_N\,R &\text{(encontrou } a \text{ depois dos } b\text{'s: rejeitar)},\\
  q_9 &= s_2\,B\,B\,s_Y\,R &\text{(chegou ao fim com pelo menos um } b\text{: aceitar)}.
\end{align*}

\medskip
\noindent\textbf{Exemplo de execução} com $W = abb$:
\[
  s_0\,abb \xrightarrow{q_3} a\,s_1\,bb \xrightarrow{q_6} ab\,s_2\,b
  \xrightarrow{q_7} abb\,s_2\,B \xrightarrow{q_9} \text{ACEITAR.}
\]

\newpage
%% ═════════════════════════════════════════════════════════════════════════════
\part*{Capítulo 14 — Conjuntos Ordenados e Reticulados}
\addcontentsline{toc}{part}{Capítulo 14 — Conjuntos Ordenados e Reticulados}

\begin{teoriacaixa}
\textbf{Ordem parcial.} Uma relação $R$ em $S$ é uma \keyword{ordem parcial} se é
reflexiva [$aRa$], antissimétrica [$aRb \wedge bRa \Rightarrow a = b$] e transitiva
[$aRb \wedge bRc \Rightarrow aRc$]. O par $(S,R)$ chama-se \keyword{conjunto parcialmente
ordenado} (poset).
\end{teoriacaixa}

%% ─────────────────────────────────────────────────────────────────────────────
\section*{Exercício 14.1 (= Exemplo 14.1 do livro) — Exemplos de Ordens Parciais}
\addcontentsline{toc}{section}{Exercício 14.1}

\subsection*{(a) Inclusão de conjuntos}

Seja $S$ uma colecção de conjuntos. A relação $\subseteq$ é uma ordem parcial:
\begin{itemize}
  \item \textbf{Reflexiva:} $A \subseteq A$ para todo $A$.
  \item \textbf{Antissimétrica:} se $A \subseteq B$ e $B \subseteq A$, então $A = B$.
  \item \textbf{Transitiva:} se $A \subseteq B$ e $B \subseteq C$, então $A \subseteq C$.
\end{itemize}

\subsection*{(b) Divisibilidade em $\mathbb{N}$}

Em $\mathbb{N}$, escrevemos $a \mid b$ se $\exists c \in \mathbb{Z}: b = ac$.
A divisibilidade é uma ordem parcial em $\mathbb{N}$:
\begin{itemize}
  \item \textbf{Reflexiva:} $a \mid a$ (tome $c = 1$).
  \item \textbf{Antissimétrica:} se $a \mid b$ e $b \mid a$ com $a,b > 0$, então $a = b$.
  \item \textbf{Transitiva:} se $a \mid b$ e $b \mid c$, então $a \mid c$.
\end{itemize}

\subsection*{(c) Divisibilidade em $\mathbb{Z}$ — não é ordem parcial}

Em $\mathbb{Z}$, a divisibilidade \emph{não} é antissimétrica: $2 \mid -2$ e $-2 \mid 2$
mas $2 \ne -2$. Logo não define uma ordem parcial em $\mathbb{Z}$.

\subsection*{(d) Relação de potência em $\mathbb{Z}$}

Definindo $aRb$ se $\exists r \in \mathbb{N}: b = a^r$. Esta relação é uma ordem parcial:
\begin{itemize}
  \item \textbf{Reflexiva:} $aRa$ pois $a = a^1$.
  \item \textbf{Antissimétrica:} se $b = a^r$ e $a = b^s$, então $b = (b^s)^r = b^{sr}$,
        logo $sr = 1$, logo $s = r = 1$ e $a = b$.
  \item \textbf{Transitiva:} se $b = a^r$ e $c = b^s$, então $c = (a^r)^s = a^{rs}$,
        logo $aRc$.
\end{itemize}

%% ─────────────────────────────────────────────────────────────────────────────
\section*{Exercício 14.2 (= Exemplo 14.2) — Comparabilidade e Ordem Linear}
\addcontentsline{toc}{section}{Exercício 14.2}

\subsection*{(a) Divisibilidade em $\mathbb{N}$}

Dois elementos $a, b$ são \keyword{comparáveis} se $a \mid b$ ou $b \mid a$.
Por exemplo, $7 \mid 21$ logo $7$ e $21$ são comparáveis. Mas $3 \nmid 5$ e $5 \nmid 3$,
logo 3 e 5 são \keyword{não-comparáveis}.

$\mathbb{N}$ com divisibilidade \emph{não} é linearmente ordenado. Porém o subconjunto
$A = \{2, 6, 12, 36\}$ é linearmente ordenado: $2 \mid 6 \mid 12 \mid 36$.

\subsection*{(b) Ordem usual em $\mathbb{N}$}

Com $\le$, quaisquer dois inteiros positivos são comparáveis, logo $\mathbb{N}$ é
linearmente ordenado, e todo o subconjunto é também linearmente ordenado.

\subsection*{(c) Inclusão em $\mathcal{P}(A)$}

Se $|A| \ge 2$, existem $\{a\}$ e $\{b\}$ com $a \ne b$ que são não-comparáveis por $\subseteq$.
Logo $\mathcal{P}(A)$ não é linearmente ordenado.
Contudo, $\{\emptyset, \{a\}, A\}$ é um subconjunto linearmente ordenado.

%% ─────────────────────────────────────────────────────────────────────────────
\section*{Exercícios 14.3--14.5 (= Exemplos 14.3, 14.4, 14.5) — Diagramas de Hasse}
\addcontentsline{toc}{section}{Exercícios 14.3--14.5}

\begin{teoriacaixa}
\textbf{Diagrama de Hasse.} Para um poset finito $S$, o \keyword{diagrama de Hasse}
representa os pares $a \lll b$ (``$a$ é predecessor imediato de $b$''): $a < b$ e
não existe $c$ com $a < c < b$. No diagrama, $b$ fica acima de $a$ e liga-se por
uma linha. Não se representam as relações transitivas nem reflexivas (redundantes).
\end{teoriacaixa}

\subsection*{Exercício 14.3 — Três exemplos de diagramas de Hasse}

\begin{enumerate}[label=(\alph*)]
  \item $A = \{1,2,3,4,6,8,9,12,18,24\}$ ordenado por divisibilidade.
        Os pares de cobertura incluem, por exemplo, $1 \lll 2$, $1 \lll 3$,
        $2 \lll 4$, $2 \lll 6$, $3 \lll 6$, $3 \lll 9$, $4 \lll 8$, $4 \lll 12$,
        $6 \lll 12$, $6 \lll 18$, $8 \lll 24$, $12 \lll 24$, $18 \lll \text{---}$
        (18 não divide 24). O diagrama tem 1 na base e 18, 24 no topo.

  \item $B = \{a,b,c,d,e\}$ com o diagrama dado: $d \le b$, $d \le a$, $e \le c$,
        $b \le a$, $c \le a$ (estes são os pares de cobertura: $d \lll b$, $e \lll c$,
        $b \lll a$, $c \lll a$).

  \item Uma \keyword{cadeia} finita (conjunto totalmente ordenado) tem diagrama linear
        — apenas um caminho. Por exemplo, uma cadeia de 5 elementos:
        $x \lll y \lll z \lll u \lll v$.
\end{enumerate}

\subsection*{Exercício 14.4 — Partições de $m = 5$}

As partições do inteiro 5 ordenam-se da seguinte forma: $P_1 \preceq P_2$ se os inteiros
de $P_1$ podem ser \emph{agrupados} para obter os de $P_2$ (ou, equivalentemente,
os de $P_2$ podem ser subdivididos para dar os de $P_1$).

Exemplo: $2{-}2{-}1 \preceq 3{-}2$ porque $2 + 1 = 3$. Os sete níveis do diagrama
(de baixo para cima) são:
\[
  1{-}1{-}1{-}1{-}1 \;\lll\; 2{-}1{-}1{-}1 \;\lll\; 2{-}2{-}1 \;\lll\;
  \begin{cases} 3{-}1{-}1 \\ 2{-}2{-}1 \end{cases} \;\lll\;
  \begin{cases} 4{-}1 \\ 3{-}2 \end{cases} \;\lll\; 5.
\]
A partição 5 é o elemento máximo (e último), e $1{-}1{-}1{-}1{-}1$ é o elemento mínimo
(e primeiro).

\subsection*{Exercício 14.5 — Elementos Mínimos, Máximos, Primeiros e Últimos}

Recorrendo aos diagramas do Exercício 14.3:

\begin{itemize}
  \item[(a.i)] $A$ (divisibilidade): elemento mínimo (e \emph{primeiro}) $= 1$;
        dois elementos máximos $= 18$ e $24$ (nenhum é \emph{último}).
  \item[(a.ii)] $B$: dois elementos mínimos $d$ e $e$ (nenhum é \emph{primeiro});
        único elemento máximo (e \emph{último}) $= a$.
  \item[(a.iii)] A cadeia: único mínimo e \emph{primeiro} $= x$;
        único máximo e \emph{último} $= v$.
\end{itemize}

\begin{enumerate}[label=(\alph*),start=2]
  \item $\mathcal{P}(A)$ com inclusão: $\emptyset$ é o \emph{primeiro} elemento
        ($\emptyset \subseteq X$ para todo $X$); $A$ é o \emph{último} elemento
        ($Y \subseteq A$ para todo $Y \in \mathcal{P}(A)$).
\end{enumerate}

%% ─────────────────────────────────────────────────────────────────────────────
\section*{Exercícios 14.6--14.30 — Supremo, Ínfimo e Reticulados}
\addcontentsline{toc}{section}{Exercícios 14.6--14.30}

\begin{teoriacaixa}
\textbf{Supremo e ínfimo.} Dado $A \subseteq S$ (poset), um \keyword{majorante} de $A$
é um $M \in S$ com $x \preceq M$ para todo $x \in A$. O \keyword{supremo} $\sup(A)$
é o menor majorante (se existir). Analogamente, um \keyword{minorante} de $A$ é $m \in S$
com $m \preceq y$ para todo $y \in A$; o \keyword{ínfimo} $\inf(A)$ é o maior minorante.

\textbf{Reticulado.} Um poset $L$ é um \keyword{reticulado} se para todo par $a,b \in L$
existem $\sup(a,b)$ e $\inf(a,b)$. Usam-se as notações $a \vee b = \sup(a,b)$ e
$a \wedge b = \inf(a,b)$.
\end{teoriacaixa}

\subsection*{Exercício 14.6 (= Exemplo 14.6)}

\begin{enumerate}[label=(\alph*)]
  \item $S = \{a,b,c,d,e,f\}$, $A = \{b,c,d\}$.
        Majorantes de $A$: apenas $e$ e $f$ (ambos sucedem $b$, $c$ e $d$).
        Como $e$ e $f$ são não-comparáveis, $\sup(A)$ \emph{não existe}.
        Minorantes de $A$: $a$ e $b$ (precedem $b$, $c$ e $d$). Mas $b \succ a$,
        logo $b$ é o maior minorante: $\inf(A) = b$.

  \item $S = \{1,\ldots,8\}$, $A = \{4,5,7\}$.
        Majorantes: 1, 2, 3 (no diagrama, todos sucedem 4, 5 e 7). Como $3 \preceq 1$
        e $3 \preceq 2$, o menor majorante é $\sup(A) = 3$.
        Minorantes: apenas $8$ (o único que precede 4, 5 e 7). Logo $\inf(A) = 8$.
\end{enumerate}

\subsection*{Exercício 14.7 (= Exemplo 14.7) — $\gcd$ e $\text{lcm}$ como $\inf$ e $\sup$}

\begin{enumerate}[label=(\alph*)]
  \item Em $\mathbb{N}$ ordenado por divisibilidade, para quaisquer $a, b$:
        \begin{itemize}
          \item Todo o divisor comum de $a$ e $b$ divide $\gcd(a,b)$, portanto
                $\gcd(a,b)$ é o maior minorante: $\inf(a,b) = \gcd(a,b)$.
          \item $\text{lcm}(a,b)$ divide todo o múltiplo comum de $a$ e $b$, portanto
                $\text{lcm}(a,b)$ é o menor majorante: $\sup(a,b) = \text{lcm}(a,b)$.
        \end{itemize}

  \item O conjunto $D_{36} = \{1,2,3,4,6,9,12,18,36\}$ (divisores de 36) é um reticulado
        com as mesmas operações. Por exemplo: $\gcd(4,6) = 2 = \inf(4,6)$ e
        $\text{lcm}(4,6) = 12 = \sup(4,6)$.
\end{enumerate}

\medskip
\noindent\textbf{Exercícios 14.8--14.30.} Os exercícios adicionais do capítulo aprofundam
os temas de enumerações consistentes, reticulados booleanos, reticulados distributivos,
propriedades algébricas de $\vee$ e $\wedge$, e isomorfismos de reticulados. A estrutura
de resolução segue os mesmos princípios: identificar a ordem, determinar coberturas,
calcular sup e inf, e verificar as propriedades de reticulado.

\newpage
%% ═════════════════════════════════════════════════════════════════════════════
\part*{Capítulo 15 — Álgebra de Boole}
\addcontentsline{toc}{part}{Capítulo 15 — Álgebra de Boole}

\begin{teoriacaixa}
\textbf{Álgebra de Boole.} Uma \keyword{álgebra de Boole} é um reticulado $(B, \vee, \wedge)$
complementado e distributivo. Os operadores são ``soma'' ($+$ ou $\vee$), ``produto''
($\cdot$ ou $\wedge$) e complemento ($'$). As leis de De Morgan afirmam:
$(a + b)' = a'b'$ e $(ab)' = a' + b'$.
\end{teoriacaixa}

%% ─────────────────────────────────────────────────────────────────────────────
\section*{Exercício 15.1 (= Exemplo 15.4) — Forma Soma-de-Produtos}
\addcontentsline{toc}{section}{Exercício 15.1 (15.4)}

\begin{teoriacaixa}
\textbf{Definições.}
Um \keyword{produto fundamental} é um produto de literais sem repetição. Uma
\keyword{expressão soma-de-produtos} é um produto fundamental ou uma soma de produtos
fundamentais nenhum dos quais está contido noutro (lei da absorção: $P + PQ = P$).
\end{teoriacaixa}

\textbf{Enunciado.}
\[
  E_1 = xz' + y'z + xyz' \qquad E_2 = xz' + x'yz' + xy'z.
\]
Mostre que $E_1$ não é uma expressão soma-de-produtos e simplifique-a. Verifique que
$E_2$ já é uma expressão soma-de-produtos.

\medskip
\textbf{Resolução.}

Para $E_1$: o produto $xz'$ está \emph{contido} em $xyz'$ (pois $xyz' = xz' \cdot y$
é mais restritivo). Pela lei da absorção ($P + PQ = P$):
\begin{align*}
  E_1 &= xz' + y'z + xyz'\\
      &= xz' + xyz' + y'z \quad \text{(reordenação)}\\
      &= xz'(1 + y) + y'z \quad \text{(factorização)}\\
      &= xz' + y'z. \quad \text{(lei da identidade: } 1 + y = 1 \text{)}
\end{align*}
Agora $xz'$ e $y'z$ não se contêm mutuamente ($xz'$ tem $x$ mas não $y$; $y'z$ tem $z$
mas não $x$). Logo $xz' + y'z$ é uma expressão soma-de-produtos.

Para $E_2$: verificar que nenhum dos três produtos fundamentais contém outro:
$xz'$, $x'yz'$, $xy'z$ têm variáveis e complementos distintos tal que nenhum é
submúltiplo de outro. Logo $E_2$ já é uma expressão soma-de-produtos.

%% ─────────────────────────────────────────────────────────────────────────────
\section*{Exercício 15.2 (= Exemplo 15.5) — Algoritmo para Forma Soma-de-Produtos}
\addcontentsline{toc}{section}{Exercício 15.2 (15.5)}

\textbf{Enunciado.} Reduza $E = ((xy)'z)'((x'+z)(y'+z'))'$ à forma soma-de-produtos.

\medskip
\textbf{Resolução.}

\passo{1} Aplicar De Morgan e involução para mover os complementos para dentro, até
que se apliquem apenas a variáveis:
\begin{align*}
  E &= ((xy)'z)' \cdot ((x'+z)(y'+z'))'\\
    &= \bigl((xy)'' + z'\bigr) \cdot \bigl((x'+z)' + (y'+z')'\bigr)\\
    &= (xy + z') \cdot (xz' + yz).
\end{align*}

\passo{2} Distribuir para obter soma de produtos:
\begin{align*}
  E &= xy \cdot xz' + xy \cdot yz + xz' \cdot z' + yz \cdot z'\\
    &= xyxz' + xyyz + xz'z' + yzz'.
\end{align*}

\passo{3} Simplificar usando idempotência ($xx = x$), complemento ($zz' = 0$) e
identidade ($z'z' = z'$):
\begin{align*}
  xyxz' &= xyz', &  xyyz &= xyz, &  xz'z' &= xz', & yzz' &= 0.
\end{align*}
Logo $E = xyz' + xyz + xz' + 0 = xyz' + xyz + xz'$.

\passo{4} Absorção: $xz'$ está contido em $xyz'$ (pois $xyz' = xz' \cdot y$), logo
$xz' + xyz' = xz'$. Eliminando $xyz'$ e o $0$:
\[
  E = xyz + xz'.
\]
Esta é a expressão soma-de-produtos final.

%% ─────────────────────────────────────────────────────────────────────────────
\section*{Exercício 15.3 (= Exemplo 15.6) — Forma Soma-de-Produtos Completa}
\addcontentsline{toc}{section}{Exercício 15.3 (15.6)}

\begin{teoriacaixa}
\textbf{Forma soma-de-produtos completa (canónica).} Uma expressão é \keyword{completa}
se cada produto fundamental (minterm) envolve \emph{todas} as $n$ variáveis.
Para completar um produto $P$ que não contenha a variável $x_i$, multiplica-se pelo
factor $(x_i + x_i') = 1$:
\[
  P = P(x_i + x_i') = Px_i + Px_i'.
\]
\end{teoriacaixa}

\textbf{Enunciado.} Exprima $E(x,y,z) = x(y'z)'$ na forma soma-de-produtos completa.

\medskip
\textbf{Resolução.}

\passo{1} Reduzir a uma forma soma-de-produtos (Algoritmo 15.1):
\begin{align*}
  E &= x(y'z)' = x(y'' + z') = x(y + z') = xy + xz'.
\end{align*}

\passo{2} Completar cada produto fundamental (Algoritmo 15.2). Cada produto deve conter
$x$, $y$ e $z$.

\begin{itemize}
  \item $xy$ não contém $z$: multiplicar por $(z + z')$:
        \[
          xy = xy(z + z') = xyz + xyz'.
        \]
  \item $xz'$ não contém $y$: multiplicar por $(y + y')$:
        \[
          xz' = xz'(y + y') = xyz' + xy'z'.
        \]
\end{itemize}

Reunindo (eliminando duplicados, pois $xyz'$ aparece duas vezes):
\[
  E = xyz + xyz' + xy'z'.
\]
Esta é a forma soma-de-produtos completa de $E$.

%% ─────────────────────────────────────────────────────────────────────────────
\section*{Exercício 15.4 (= Exemplo 15.7) — Consenso de Produtos Fundamentais}
\addcontentsline{toc}{section}{Exercício 15.4 (15.7)}

\begin{teoriacaixa}
\textbf{Consenso.} O \keyword{consenso} de $P_1$ e $P_2$ (quando exactamente uma
variável $x_k$ aparece não-complementada num e complementada no outro) obtém-se:
eliminando $x_k$ e $x_k'$ e multiplicando os literais restantes (sem repetição).
Lema: $P_1 + P_2 + Q = P_1 + P_2$ onde $Q$ é o consenso.
\end{teoriacaixa}

\textbf{Determine o consenso $Q$ para:}

\begin{enumerate}[label=(\alph*)]
  \item $P_1 = xyz's$ e $P_2 = xy't$: a variável $y$ aparece sem complemento em $P_1$
        e complementada em $P_2$. Elimina-se $y$ e $y'$ e multiplica-se o resto:
        \[
          Q = x \cdot z' \cdot x \cdot t = xz'st \quad \text{(sem repetição de } x\text{).}
        \]

  \item $P_1 = xy'$ e $P_2 = y$: eliminar $y$ e $y'$, restam $x$ (de $P_1$) e nada
        (de $P_2$). Logo $Q = x$.

  \item $P_1 = x'yz$ e $P_2 = x'yt$: \emph{nenhuma} variável aparece em complementos
        opostos. Não existe consenso.

  \item $P_1 = x'yz$ e $P_2 = xyz'$: \emph{duas} variáveis ($x$ e $z$) aparecem em
        complementos opostos. O consenso não é definido quando há mais de uma tal variável.
\end{enumerate}

%% ─────────────────────────────────────────────────────────────────────────────
\section*{Exercício 15.5 (= Exemplo 15.8) — Método do Consenso para Implicantes Primos}
\addcontentsline{toc}{section}{Exercício 15.5 (15.8)}

\textbf{Enunciado.} Seja $E = xyz + x'z' + xyz' + x'y'z + x'yz'$. Encontre os
implicantes primos de $E$ pelo método do consenso.

\medskip
\textbf{Resolução.}

\passo{1} Eliminar produtos redundantes (lei da absorção: se $P_i \supseteq P_j$,
eliminar $P_i$):
\[
  E = xyz + x'z' + xyz' + x'y'z \quad\text{($x'yz'$ inclui $x'z'$ pois $x'z' \subseteq x'yz'$? Não — é ao contrário: $x'z'$ está em $x'yz'$, logo elimina-se $x'yz'$).}
\]
Verificando: $x'z' = x' \cdot z'$ sem restrição em $y$, enquanto $x'yz' = x' \cdot y \cdot z'$.
Como $x'z'$ é mais geral (``inclui'' $x'yz'$ no sentido de $P+PQ=P$), elimina-se $x'yz'$:
\[
  E = xyz + x'z' + xyz' + x'y'z.
\]

\passo{2} Adicionar consensos novos. Consenso de $xyz$ e $xyz'$: variável $z$ oposta.
$Q = xy$. Adicionamos $xy$:
\[
  E = xyz + x'z' + xyz' + x'y'z + xy.
\]

\passo{3} $xyz$ e $xyz'$ estão ambos contidos em $xy$ (absorção): eliminam-se.
\[
  E = x'z' + x'y'z + xy.
\]

\passo{4} Consenso de $x'z'$ e $x'y'z$: variável $z$ oposta. $Q = x'y'$.
\[
  E = x'z' + x'y'z + xy + x'y'.
\]

\passo{5} $x'y'z$ inclui $x'y'$: eliminar $x'y'z$.
\[
  E = x'z' + xy + x'y'.
\]

\passo{6} Consenso de $x'z'$ e $xy$: variável $x$ oposta. $Q = yz'$.
\[
  E = x'z' + xy + x'y' + yz'.
\]

Nenhuma nova absorção ou consenso é possível. Os \keyword{implicantes primos} de $E$
são $x'z'$, $xy$, $x'y'$ e $yz'$.

%% ─────────────────────────────────────────────────────────────────────────────
\section*{Exercício 15.6 (= Exemplo 15.9) — Forma Mínima Soma-de-Produtos}
\addcontentsline{toc}{section}{Exercício 15.6 (15.9)}

\textbf{Enunciado.} Use os implicantes primos do Exercício 15.5 para encontrar a forma
mínima soma-de-produtos de $E = x'z' + xy + x'y' + yz'$.

\medskip
\textbf{Resolução.}

\passo{1} Expandir cada implicante primo à forma soma-de-produtos completa:
\begin{align*}
  x'z' &= x'z'(y+y') = x'yz' + x'y'z',\\
  xy   &= xy(z+z') = xyz + xyz',\\
  x'y' &= x'y'(z+z') = x'y'z + x'y'z',\\
  yz'  &= yz'(x+x') = xyz' + x'yz'.
\end{align*}

\passo{2} Os minterms de $x'z'$ são $\{x'yz',\, x'y'z'\}$. Verificar se estes aparecem
nos minterms de outros implicantes: $x'yz' \in yz'$ e $x'y'z' \in x'y'$. Ambos estão
cobertos. Logo $x'z'$ é \emph{supérfluo} e pode ser eliminado:
\[
  E = xy + x'y' + yz'.
\]

\passo{3} Verificar se $xy$, $x'y'$ ou $yz'$ são supérfluos: nenhum dos seus minterms
está coberto apenas pelos outros dois. Logo a forma mínima soma-de-produtos é:
\[
  \boxed{E = xy + x'y' + yz'.}
\]

%% ─────────────────────────────────────────────────────────────────────────────
\section*{Exercícios 15.7--15.13 — Portas Lógicas e Circuitos}
\addcontentsline{toc}{section}{Exercícios 15.7--15.13}

\begin{teoriacaixa}
\textbf{Portas lógicas básicas.}
\begin{itemize}[noitemsep]
  \item \textbf{Porta OR}: $Y = A + B$ ($Y = 0$ s.s.e. $A = B = 0$).
  \item \textbf{Porta AND}: $Y = A \cdot B$ ($Y = 1$ s.s.e. $A = B = 1$).
  \item \textbf{Porta NOT} (inversor): $Y = A'$ ($Y = 1$ s.s.e. $A = 0$).
\end{itemize}
Qualquer expressão Booleana pode ser implementada com combinações destas três portas.
Um \keyword{circuito lógico} é uma rede dessas portas.
\end{teoriacaixa}

\subsection*{Exercício 15.7 — Porta OR}

A porta OR com entradas $A$ e $B$ tem saída $Y = A + B$. A tabela de verdade:

\begin{center}
\begin{tabular}{cc|c}
  \toprule
  $A$ & $B$ & $Y = A+B$\\
  \midrule
  0 & 0 & 0\\
  0 & 1 & 1\\
  1 & 0 & 1\\
  1 & 1 & 1\\
  \bottomrule
\end{tabular}
\end{center}

A saída é 0 apenas quando \emph{ambas} as entradas são 0. Com $n$ entradas,
$Y = 0$ s.s.e. todas as entradas são 0.

\subsection*{Exercício 15.8 — Porta AND}

A porta AND com entradas $A$ e $B$ tem saída $Y = A \cdot B$. A tabela de verdade:

\begin{center}
\begin{tabular}{cc|c}
  \toprule
  $A$ & $B$ & $Y = AB$\\
  \midrule
  0 & 0 & 0\\
  0 & 1 & 0\\
  1 & 0 & 0\\
  1 & 1 & 1\\
  \bottomrule
\end{tabular}
\end{center}

A saída é 1 apenas quando \emph{ambas} as entradas são 1.

\subsection*{Exercício 15.9 — Porta NOT e Implementação de Circuitos}

A porta NOT tem $Y = A'$: inverte o bit de entrada.
Um circuito lógico que implementa $E = xy + x'y' + yz'$ (forma mínima do Exercício 15.6)
precisaria de:
\begin{enumerate}
  \item Inversores para obter $x'$, $y'$, $z'$.
  \item Três portas AND de dois inputs: $x \cdot y$, $x' \cdot y'$, $y \cdot z'$.
  \item Uma porta OR de três inputs: $xy + x'y' + yz'$.
\end{enumerate}
No total: 3 inversores + 3 AND + 1 OR = 7 portas.

\subsection*{Exercícios 15.10--15.13}

Estes exercícios aprofundam a análise e síntese de circuitos lógicos. A metodologia geral é:
\begin{enumerate}
  \item Partir da expressão Booleana (ou tabela de verdade).
  \item Minimizar usando soma-de-produtos e implicantes primos.
  \item Implementar com portas OR, AND e NOT.
  \item Verificar equivalência através da tabela de verdade.
\end{enumerate}

\newpage
%% ═════════════════════════════════════════════════════════════════════════════
\part*{Apêndice B — Sistemas Algébricos}
\addcontentsline{toc}{part}{Apêndice B — Sistemas Algébricos}

\begin{teoriacaixa}
\textbf{Conceitos fundamentais.}
Dado um conjunto $S$ com uma operação binária $*$:
\begin{itemize}[noitemsep]
  \item $*$ é \keyword{associativa} se $(a * b) * c = a * (b * c)$.
  \item $*$ é \keyword{comutativa} se $a * b = b * a$.
  \item $e$ é \keyword{elemento neutro} se $a * e = e * a = a$.
  \item $b$ é o \keyword{inverso} de $a$ se $a * b = b * a = e$.
  \item $(S, *)$ é um \keyword{semigrupo} se $*$ é associativa.
  \item $(S, *)$ é um \keyword{monóide} se é semigrupo com elemento neutro.
  \item $(S, *)$ é um \keyword{grupo} se é monóide e todo elemento tem inverso.
\end{itemize}
\end{teoriacaixa}

%% ─────────────────────────────────────────────────────────────────────────────
\section*{Exercício B.1 (= Exemplo B.4) — Inversos em $\mathbb{Q}$}
\addcontentsline{toc}{section}{Exercício B.1 (B.4)}

\textbf{Enunciado.} Nos racionais $\mathbb{Q}$, identifique os elementos neutros e os
inversos para a adição e a multiplicação.

\medskip
\textbf{Resolução.}

\begin{enumerate}[label=(\alph*)]
  \item \textbf{Adição:} elemento neutro $= 0$ (pois $a + 0 = 0 + a = a$).
        O inverso de $a$ é $-a$ (pois $a + (-a) = 0$).
        Exemplo: os inversos aditivos de $-3$ e $3$ são $3$ e $-3$, respectivamente,
        pois $(-3) + 3 = 3 + (-3) = 0$.

  \item \textbf{Multiplicação:} elemento neutro $= 1$ (pois $a \cdot 1 = 1 \cdot a = a$).
        O inverso de $a \ne 0$ é $1/a$ (pois $a \cdot (1/a) = 1$).
        Exemplo: os inversos multiplicativos de $-3$ são $-1/3$:
        $(-3)(-1/3) = (-1/3)(-3) = 1$.
        \textbf{Atenção:} $0$ não tem inverso multiplicativo (divisão por zero indefinida).
\end{enumerate}

\medskip
\noindent\textit{Nota estrutural:} $(\mathbb{Q}, +)$ é um grupo (abeliano). $(\mathbb{Q}\setminus\{0\}, \cdot)$
é também um grupo (abeliano). Em conjunto, $\mathbb{Q}$ é um \keyword{corpo}.

%% ─────────────────────────────────────────────────────────────────────────────
\section*{Exercício B.2 (= Exemplo B.5) — Semigrupos e Monóides}
\addcontentsline{toc}{section}{Exercício B.2 (B.5)}

\textbf{Enunciado.} Classifique como semigrupo ou monóide:
$(a)$ $(\mathbb{N}, +)$ e $(\mathbb{N}, \times)$;\;
$(b)$ $F(S)$ (funções $S \to S$ com composição);\;
$(c)$ $(S, *)$ e $(S, \cdot)$ com as tabelas dadas.

\medskip
\textbf{Resolução.}

\begin{enumerate}[label=(\alph*)]
  \item $(\mathbb{N}, +)$: adição é associativa, logo é semigrupo. Mas $\mathbb{N}$ não
        contém $0$ (inteiros positivos), pelo que \emph{não} há elemento neutro.
        \textbf{Semigrupo, não monóide.}

        $(\mathbb{N}, \times)$: multiplicação é associativa e $1 \in \mathbb{N}$ é
        elemento neutro. \textbf{Monóide.}

  \item $F(S)$ (composição $\circ$): composição de funções é sempre associativa ($(f \circ g) \circ h = f \circ (g \circ h)$). A função identidade $\text{id}_S$ serve como elemento neutro. \textbf{Monóide.}

  \item Para a operação $*$ definida por $x * y = x$ (toma sempre o valor do primeiro operando):
        \[
          (x * y) * z = x * z = x \quad\text{e}\quad x * (y * z) = x * y = x.
        \]
        Logo $*$ é associativa. \textbf{Semigrupo} (mas verificar se há neutro: $e * y = e$
        para todo $y$ pede $e = e$, e $y * e = y$ pede $y = y$; portanto qualquer elemento
        é neutro à direita, mas nenhum é neutro à esquerda se o conjunto tiver mais de um
        elemento, logo não é monóide).

        Para a operação $\cdot$ com a tabela dada: verificou-se que
        $(b \cdot c) \cdot c = a \cdot c = c$ mas $b \cdot (c \cdot c) = b \cdot a = b$.
        Como $c \ne b$, $\cdot$ \textbf{não é associativa}, logo $(S, \cdot)$ \textbf{não é semigrupo}.
\end{enumerate}

%% ─────────────────────────────────────────────────────────────────────────────
\section*{Exercício B.3 (= Exemplo B.6) — Subsemigrupos}
\addcontentsline{toc}{section}{Exercício B.3 (B.6)}

\textbf{Enunciado.} Determine quais dos subconjuntos são subsemigrupos.

\begin{enumerate}[label=(\alph*)]
  \item $A$ (pares) e $B$ (ímpares) em $(\mathbb{N}, \times)$ e $(\mathbb{N}, +)$.

  \medskip
  \textbf{Resolução.}
  \begin{itemize}
    \item $(A, \times)$: par $\times$ par $=$ par. $A$ é fechado. \textbf{Subsemigrupo.}
    \item $(B, \times)$: ímpar $\times$ ímpar $=$ ímpar. $B$ é fechado. \textbf{Subsemigrupo.}
    \item $(A, +)$: par $+$ par $=$ par. $A$ é fechado. \textbf{Subsemigrupo.}
    \item $(B, +)$: ímpar $+$ ímpar $=$ par $\notin B$. $B$ \textbf{não} é fechado.
          \textbf{Não é subsemigrupo.}
  \end{itemize}

  \item $H$ (palavras de comprimento par) no semigrupo livre $F = F(\{a,b\})$.

  \medskip
  \textbf{Resolução.} Se $u, v \in H$ têm comprimento par, então $|uv| = |u| + |v|$
  é par $+$ par $=$ par. Logo $uv \in H$. $H$ é fechado. \textbf{Subsemigrupo.}
\end{enumerate}

%% ─────────────────────────────────────────────────────────────────────────────
\section*{Exercício B.4 (= Exemplo B.7) — Relações de Congruência}
\addcontentsline{toc}{section}{Exercício B.4 (B.7)}

\begin{teoriacaixa}
\textbf{Relação de congruência.} Uma relação de equivalência $\sim$ em $S$ é uma
\keyword{relação de congruência} se for compatível com a operação:
$a \sim a'$ e $b \sim b'$ implica $ab \sim a'b'$.
O \keyword{quociente} $S/{\sim}$ é então um semigrupo com $[a][b] = [ab]$.
\end{teoriacaixa}

\textbf{Enunciado.} Verifique que:
$(a)$ em $F(\{a,b,\ldots\})$, a relação $u \sim u'$ se $|u| = |u'|$ é congruência;
$(b)$ a congruência módulo $m$ em $\mathbb{Z}$ é uma relação de congruência.

\medskip
\textbf{Resolução.}

\begin{enumerate}[label=(\alph*)]
  \item Se $|u| = |u'| = p$ e $|v| = |v'| = q$, então:
        \[
          |uv| = |u| + |v| = p + q = |u'| + |v'| = |u'v'|.
        \]
        Logo $uv \sim u'v'$. É uma relação de congruência.

        O quociente $F/{\sim}$ identifica-se com $(\mathbb{N}, +)$: cada classe de
        equivalência é identificada pelo comprimento da palavra.

  \item Em $\mathbb{Z}$, se $a \equiv c \pmod{m}$ e $b \equiv d \pmod{m}$, então:
        \begin{align*}
          a + b &\equiv c + d \pmod{m},\\
          ab    &\equiv cd \pmod{m}.
        \end{align*}
        Estas propriedades (Teorema 11.22 do livro) mostram que $\equiv \pmod{m}$ é
        compatível com adição e multiplicação, sendo portanto uma relação de congruência
        em $(\mathbb{Z}, +)$ e em $(\mathbb{Z}, \cdot)$.

        O quociente $\mathbb{Z}/{\equiv_m}$ é o conjunto $\mathbb{Z}_m = \{0, 1, \ldots, m-1\}$
        com aritmética módulo $m$.
\end{enumerate}

%% ─────────────────────────────────────────────────────────────────────────────
\section*{Exercício B.5 — Unicidade do Elemento Neutro e dos Inversos}
\addcontentsline{toc}{section}{Exercício B.5}

\textbf{Enunciado (Teorema B.2).} Se $e$ é neutro à esquerda e $f$ é neutro à direita
para uma operação $*$ em $S$, então $e = f$.

\medskip
\textbf{Resolução.}

\begin{align*}
  e * f &= f \quad\text{(porque } e \text{ é neutro à esquerda)},\\
  e * f &= e \quad\text{(porque } f \text{ é neutro à direita)}.
\end{align*}
Logo $e = f$. $\square$

\medskip
\noindent\textbf{Corolário.} O elemento neutro (se existir) é \emph{único}. De facto,
se $e_1$ e $e_2$ são ambos neutros, então $e_1 = e_1 * e_2 = e_2$.

%% ─────────────────────────────────────────────────────────────────────────────
\section*{Exercício B.6 — Leis de Cancelamento}
\addcontentsline{toc}{section}{Exercício B.6}

\textbf{Enunciado.} A operação de multiplicação de matrizes não satisfaz as leis de
cancelamento. Verifique com o exemplo dado.

\medskip
\textbf{Resolução.}

No livro, os seguintes dados são fornecidos:
\[
  A = \begin{pmatrix} 1 & 1 \\ 0 & 0 \end{pmatrix}, \quad
  B = \begin{pmatrix} 1 & 1 \\ 0 & 1 \end{pmatrix}, \quad
  C = \begin{pmatrix} 0 & -3 \\ 1 & 5 \end{pmatrix}, \quad
  D = \begin{pmatrix} 1 & 2 \\ 0 & 0 \end{pmatrix}.
\]

Calculemos $AB$ e $AC$:
\begin{align*}
  AB &= \begin{pmatrix}1&1\\0&0\end{pmatrix}\begin{pmatrix}1&1\\0&1\end{pmatrix}
       = \begin{pmatrix}1\cdot1+1\cdot0 & 1\cdot1+1\cdot1\\
                        0\cdot1+0\cdot0 & 0\cdot1+0\cdot1\end{pmatrix}
       = \begin{pmatrix}1&2\\0&0\end{pmatrix} = D,\\[6pt]
  AC &= \begin{pmatrix}1&1\\0&0\end{pmatrix}\begin{pmatrix}0&-3\\1&5\end{pmatrix}
       = \begin{pmatrix}0+1 & -3+5\\0+0 & 0+0\end{pmatrix}
       = \begin{pmatrix}1&2\\0&0\end{pmatrix} = D.
\end{align*}

Portanto $AB = AC = D$, mas $B \ne C$. Isto demonstra que a lei de cancelamento
à esquerda \emph{falha} para a multiplicação de matrizes: $AB = AC$ \emph{não}
implica $B = C$.

\medskip
\noindent\textit{Porquê?} A lei de cancelamento é equivalente a dizer que a
matriz $A$ é invertível ($\det A \ne 0$). Neste caso,
$\det A = 1 \cdot 0 - 1 \cdot 0 = 0$, logo $A$ é singular e não cancelável.

%% ═════════════════════════════════════════════════════════════════════════════
\newpage
\section*{Conclusão}
\addcontentsline{toc}{section}{Conclusão}

Este documento cobriu os exercícios resolvidos dos seguintes capítulos:

\begin{itemize}
  \item \textbf{Capítulo 10} (Exercícios 10.10--10.12): estrutura e representação de
        árvores binárias (árvores similares, representação ligada e sequencial).
  \item \textbf{Capítulo 12} (Exercícios 12.1--12.21): linguagens formais, expressões
        regulares e autómatos finitos, incluindo construções clássicas (complementar,
        interseção, união).
  \item \textbf{Capítulo 13} (Exercícios 13.1--13.5): máquinas de Turing — quadros,
        transições, ciclos e reconhecimento de linguagens.
  \item \textbf{Capítulo 14} (Exercícios 14.1--14.30): conjuntos parcialmente ordenados,
        diagramas de Hasse, supremo, ínfimo e reticulados.
  \item \textbf{Capítulo 15} (Exercícios 15.1--15.13): álgebra de Boole, simplificação
        de expressões, forma canónica e circuitos lógicos.
  \item \textbf{Apêndice B} (Exercícios B.1--B.6): sistemas algébricos — semigrupos,
        monóides, substrutaturas, relações de congruência e leis de cancelamento.
\end{itemize}

\end{document}
